
\documentclass[10pt]{article} % For LaTeX2e
\usepackage{tmlr}
% If accepted, instead use the following line for the camera-ready submission:
%\usepackage[accepted]{tmlr}
% To de-anonymize and remove mentions to TMLR (for example for posting to preprint servers), instead use the following:
%\usepackage[preprint]{tmlr}

% Optional math commands from https://github.com/goodfeli/dlbook_notation.
\input{math_commands.tex}

\usepackage{hyperref}
\usepackage{url}
\usepackage[capitalize,noabbrev]{cleveref}

\usepackage{nicefrac}
\usepackage{bm}
\usepackage{graphicx}
\usepackage{caption}
\usepackage{subcaption}
\usepackage{booktabs}
\usepackage{multirow}

\usepackage{amsmath}
\usepackage{amssymb}
\usepackage{mathtools}
\usepackage{amsthm}

\theoremstyle{plain}
\newtheorem{theorem}{Theorem}[section]
\newtheorem{proposition}[theorem]{Proposition}
\newtheorem{lemma}[theorem]{Lemma}
\newtheorem{corollary}[theorem]{Corollary}
\theoremstyle{definition}
\newtheorem{definition}[theorem]{Definition}
\newtheorem{assumption}[theorem]{Assumption}
\theoremstyle{remark}
\newtheorem{remark}[theorem]{Remark}

\title{Wrapped Flow Matching for Symmetric Positive Definite Matrices}

% Authors must not appear in the submitted version. They should be hidden
% as long as the tmlr package is used without the [accepted] or [preprint] options.
% Non-anonymous submissions will be rejected without review.

\author{\name Thibault de Surrel \email thibault.de-surrel@lamsade.dauphine.fr \\
  \addr LAMSADE, CNRS,\\
  Université Paris Dauphine-PSL, \\
  Paris, France
  \AND
  \name Aimi Okabayashi \email aimi.okabayashi@univ-ubs.fr \\
  \addr IRISA, UMR CNRS 6074, \\
  Université Bretagne Sud,
  Vannes, France
}

% The \author macro works with any number of authors. Use \AND
% to separate the names and addresses of multiple authors.

\newcommand{\fix}{\marginpar{FIX}}
\newcommand{\new}{\marginpar{NEW}}

\def\month{MM}  % Insert correct month for camera-ready version
\def\year{YYYY} % Insert correct year for camera-ready version
\def\openreview{\url{https://openreview.net/forum?id=XXXX}} % Insert correct link to OpenReview for camera-ready version

\begin{document}

\maketitle

\begin{abstract}
  Conditional Flow Matching (CFM) has emerged as a state-of-the-art paradigm for generative modeling, yet its extension to non-Euclidean geometries remains an open challenge. Current Riemannian CFM methods rely strictly on geodesic interpolation or on pullback to Euclidean spaces, limiting their expressivity. In this work, we propose a flexible formulation for Riemannian CFM that leverages the wrapped Gaussian distribution to construct conditional probability paths. Our approach provides a unified framework that encompasses existing Riemannian CFM variants, enabling more adaptable flows that respect the intrinsic geometry of the manifold. Moreover, we are able to derive from this framework a new method that we coined \emph{WrappedCFM}. We focus our development on the manifold of Symmetric Positive Definite (SPD) matrices denoted $\P_d$, a core component of information geometry. Our geometric formulation respects the intrinsic structure of $\P_d$ and relies neither on extrinsic Euclidean approximations nor on geodesic interpolation. Through both synthetic and real-world experiments, we demonstrate that WrappedCFM accurately models complex distributions on $\P_d$ and achieves competitive performance, showcasing the advantages of flexible path construction over geodesic interpolations.

\end{abstract}

\section{Introduction}
Denoising Diffusion Probabilistic Models (DDPM) \citep{ho2020denoising, song2021scorebased} and Conditional Flow Matching (CFM) \citep{lipmanFlowMatchingGenerative2023, liu2022flow, jung2024flowavse} have established themselves as the state-of-the-art for generative modeling on Euclidean data. However, real-world data often possess a richer geometric structure, and modeling them as lying on a manifold is often a powerful solution \citep{Papillon_2025}. This assumption requires to adapt classical methods to capture the intrinsic geometry of the underlying space \citep{pennecIntrinsicStatisticsRiemannian2006, guigui2023introduction}.

\textbf{In this paper, we bridge this gap by extending CFM algorithms to the Riemannian manifold $\P_d$ of Symmetric Positive Definite (SPD) matrices.} SPD matrices are widely used throughout statistics and machine learning, appearing in domains such as \emph{Brain Computer Interfaces} (BCI) \citep{lotteReviewClassificationAlgorithms2018}, \emph{Diffusion Tensor Imaging} \citep{pennecManifoldvaluedImageProcessing2020}, \emph{radar signals} \citep{brooks} or \emph{hyperspectral imaging} \citep{russwurm2019breizhcrops}. This manifold is also of central interest in information geometry where it is associated to the statistical model of multivariate centered Gaussians \citep{skovgaardRiemannianGeometryMultivariate1984}. The space $\P_d$ is a curved Riemannian manifold, and taking into account this non-flat geometry is crucial when manipulating SPD matrices (see for example the \emph{swelling effect} in \citealt{arsigny2007geometric}).

In this paper, we propose two main contributions.
\begin{enumerate}
  \item \textbf{First, the key novelty} is a general framework for CFM on $\P_d$, encompassing state-of-the-art Riemannian CFM methods, namely \emph{RiemannianCFM} \citep{chen2024flow} and \emph{DiffeoCFM} \citep{collasRiemannianFlowMatching2025}. For this, we leverage an infinitely supported distribution directly defined on the manifold $\P_d$, the \emph{wrapped Gaussian} \citep{desurrel2025wrappedgaussianmanifoldsymmetric}, to design the conditional probability paths.
  \item \textbf{Secondly}, using the flexibility of our framework, we propose a novel interpolation scheme that extends beyond restricted geodesic paths. We call this method \emph{WrappedCFM}.
\end{enumerate}

This paper is organized as follows: after reviewing state-of-the-art CFM methods on Riemannian manifold in \cref{sec:related_works}, we introduce the Riemannian geometry of $\P_d$ as well as the wrapped Gaussian distribution in \cref{sec:spd_geom}. Next, we recall the classical theory of Euclidean CFM and review two SOTA Riemannian CFM methods in \cref{sec:background_CFM}. We then introduce our unified framework in \cref{sec:global_framework} and propose our new interpolation \emph{WrappedCFM} in \cref{sec:W-CFM}. We finally validate our method on synthetic data in \cref{sec:synthetic_exp} and on real-world data in \cref{sec:real_world_experiments}.

\section{Related works}
\label{sec:related_works}
The works of \citet{ho2020denoising, song2021scorebased} on Denoising Diffusion Probabilistic Models (DDPM) and of \citet{lipmanFlowMatchingGenerative2023, liu2022flow, jung2024flowavse} on Conditional Flow Matching (CFM) have led to state-of-the-art generative models. Several attempts have been made to generalize these models to Riemannian manifolds. In \citet{liSPDDDPMDenoisingDiffusion2023}, the authors propose a DDPM specific to the manifold $\P_d$ of SPD matrices: using the isotropic Gaussian from \citet{saidRiemannianGaussianDistributions2016}, they adapt the classical forward and backward process of DDPM to $\P_d$. Other Riemannian adaptations of the diffusion generative framework have been proposed for general Riemannian manifolds \citep{jo2023generative, de2022riemannian, huang2022riemannian} that learn to reverse the noising process. However, most of these works struggle to perform on complex manifolds (outside of the sphere or the torus) as they rely on approximations of the intractable heat kernel. Moreover, they sometimes lack a clear geometrical interpretation, which will not be the case for the framework that we propose.

Regarding CFM, \citet{chen2024flow} propose to adapt the Euclidean CFM framework to the setup of a Riemannian manifold. They generalize the linear interpolation used to build conditional probability paths to geodesic interpolation, building Riemannian conditional probability paths based on geodesics. In our work, we focus on infinity supported Gaussian conditional probability paths. The RiemannianCFM with geodesic interpolation has been used in \citet{boseSE3StochasticFlowMatching2024} to generate novel protein structures working on the special Euclidean group $\text{SE}(3)$ or in \citet{braun2024riemannian} to learn and generate robot sensorimotor policies working on the hypersphere or on $\text{SO}(3)$. In \citet{zaghen2025riemannianvariationalflowmatching}, the authors developed a Riemannian version of Variational Flow Matching, based on the isotropic Gaussian of \citet{saidRiemannianGaussianDistributions2016}. Another proposition to adapt CFM to Riemannian manifolds is by using pullback geometries, as in \citet{collasRiemannianFlowMatching2025}. Their so-called \emph{DiffeoCFM} sends the data to a Euclidean space $E$ (for example a specific tangent space) \emph{via} a global diffeomorphism and relies on Euclidean CFM in $E$ before sending back the generated samples to the manifold. Our approach is different as we work directly on the manifold without relying on any artificial Euclidean space. In their paper, the authors apply their DiffeoCFM to fMRI and EEG datasets. A similar approach is \citet{zhoubrainflow}, in which the authors use the Cholesky decomposition to perform efficient FM on the simpler manifold of lower-triangular matrices with positive diagonal. They also lead experiments on fMRI datasets.

%Parler de \cite{zaghenVariationalFlowMatching2025} ?

\section{The manifold of SPD matrices}
\label{sec:spd_geom}
\textbf{Bold uppercase letters} ($\Xm, \Ym, \Um, \Vm$) denote matrices. For $d \in \mathbb{N}^\star$, let $D := d(d+1)/2$. $\Im_d$ is the $d \times d$ identity matrix, while $\exp$ and $\log$ respectively refer to the matrix exponential and logarithm.
\subsection{The Riemannian geometry of SPD matrices}
The set of $d \times d$ \emph{Symmetric Positive Definite} (SPD) matrices, denoted by $\P_d$, is defined as
\[\P_d = \{\Xm \in \R^{d \times d} \mid \Xm^\top = \Xm \text{ and } \forall z \neq 0,~z^\top \Xm z > 0\}.\]

This set is an open convex cone in the space $\S_d$ of symmetric matrices. Therefore, it is a manifold of dimension $D = d(d+1)/2$ and the tangent space $T_{\Xm} \P_d$ at any point $\Xm \in \P_d$ is identified with $\S_d$. A Riemannian metric can be defined on $\P_d$ by endowing each tangent space $T_{\Xm}\P_d$ with the following inner product, called the \emph{Affine Invariant Riemannian Metric} (AIRM) \citep{pennecManifoldvaluedImageProcessing2020}:
\begin{equation}
  \forall \Um,\Vm \in T_{\Xm}\P_d,~\langle \Um,\Vm\rangle_{\Xm} = \tr(\Xm^{-1}\Um\Xm^{-1}\Vm).
  \label{eq:AIRM_metric}
\end{equation}
Once endowed with this metric, $\P_d$ is a geodesically complete manifold with non-positive curvature, thus $\P_d$ forms a \emph{Hadamard manifold} \citep{shigaHadamardManifolds1984}. The AIRM distance $\delta$ between two points $\Xm,\Ym \in \P_d$ is:
\begin{equation}
  \label{eq:distance_spd}
  \delta(\Xm,\Ym) = \|\log(\Xm^{-1/2}\Ym\Xm^{-1/2})\|_F,
\end{equation}
where $\|\cdot\|_F$ is the Frobenius norm. We will also need the exponential map $\Exp_{\Xm} : T_{\Xm}\P_d \to \P_d$ and its inverse, the logarithmic map $\Log_{\Xm} : \P_d \to T_{\Xm}\P_d$. As $\P_d$ is a Hadamard manifold, these maps are diffeomorphisms. Moreover, \citet{bhatiaPositiveDefiniteMatrices2007} gives them in closed form:
\begin{equation}
  \label{eq:exp_log_riem}
  \begin{aligned}
    \Exp_{\Xm}(\Um) &= \Xm^{1/2}\exp(\Xm^{-1/2}\Um\Xm^{-1/2})\Xm^{1/2}, \\
    \Log_{\Xm}(\Ym) &= \Xm^{1/2}\log(\Xm^{-1/2}\Ym\Xm^{-1/2})\Xm^{1/2}.
  \end{aligned}
\end{equation}
Finally, when $\P_d$ is equipped with the AIRM metric given at \cref{eq:AIRM_metric}, the Riemannian volume element at $\Xm = [[x_{ij}]] \in \P_d$ is given by (see Section 4.1.3 of \citealt{terrasHarmonicAnalysisSymmetric1988}):
\begin{equation}
  \label{eq:volume}
  \mathrm{d} \text{vol}(\Xm) = \det(\Xm)^{-(d+1)/2} \prod_{i \leq j} \mathrm{d}x_{ij},
\end{equation}
where $\mathrm{d}x_{ij}$ is the Lebesgue measure on $\R$.

\begin{remark}
  Other Riemannian metrics have been defined on the manifold $\P_d$ of SPD matrices such as the \emph{log-Euclidean metric} \citep{arsigny2007geometric}, the \emph{Cholesky metric} \citep{dryden2009non} or the \emph{Bures-Wasserstein metric} \citep{bhatia2019bures}. In this work, we decided to focus on the AIRM metric as it leads to a curved space (unlike the log-Euclidean nor the Cholesky metric) with a non-positive curvature (unlike the Bures-Wasserstein metric). This makes the space interesting (non-flat) with nice mathematical properties ($\Exp_{\Xm}$ is a diffeomorphism). Moreover, it is one of the most used metrics in practice \citep{chevallier2021review}.
\end{remark}

\subsection{Wrapped Gaussian on \texorpdfstring{$\P_d$}{Pd}}
\label{sec:wrapped_gaussian}
\begin{figure}[t]
  \centering
  \includegraphics[width=0.5\linewidth]{figures/wg.pdf}
  \caption{An illustration of the wrapped Gaussian $\WG(\Xm, \mu, \Sigma)$ defined on the manifold $\P_d$ of SPD matrices.}
  \label{fig:wg}
\end{figure}

There are several ways to define a Gaussian distribution on $\P_d$ \citep{pennecIntrinsicStatisticsRiemannian2006,saidRiemannianGaussianDistributions2016}. In this work, we consider the \emph{wrapped Gaussian distribution} \citep{desurrel2025wrappedgaussianmanifoldsymmetric, chevallierExponentialWrappedDistributionsSymmetric2022}, which is defined as the pushforward of a Euclidean multivariate Gaussian distribution from a tangent space $T_{\Xm}\P_d$ at a given point $\Xm \in \P_d$ by the exponential map $\Exp_{\Xm}$. In order to simplify the notations, we consider in all this paper the affine-invariant (AIRM-orthonormal) identification of $T_{\Xm}\P_d$ with $\R^D$, i.e. a tangent vector $\Um \in T_{\Xm}\P_d$ is represented by the coordinates of $\Xm^{-1/2}\Um\Xm^{-1/2}$ in a fixed Frobenius-orthonormal basis of $\mathcal{S}_d$. More on this in \citet{desurrel2025wrappedgaussianmanifoldsymmetric} or in Section 3.3.3.3 of \citet{pennecManifoldvaluedImageProcessing2020}.
\begin{definition}[Wrapped Gaussian]
  \label{def:wrapped_gaussian}
  Let $\Xm \in \P_d$, $\mu \in \R^D$ and $\Sigma \in \P_{D}$. A random variable $\Zm$ on $\P_d$ follows a wrapped Gaussian (WG) denoted $\WG(\Xm,\mu, \Sigma)$ if
  $$\Zm = \Exp_{\Xm}(\mathbf{t}),~ \mathbf{t} \sim \N(\mu, \Sigma).$$
\end{definition}
In other words, $\WG(\Xm,\mu, \Sigma) = \Exp_{\Xm}\#\N(\mu, \Sigma),$ where $\#$ denotes the push-forward. This distribution relies on three parameters: $\Xm \in \P_d$ is the anchor points of the distribution, that gives a tangent space from which the Euclidean Gaussian will be push-forward. Then, $\mu \in \R^D$ is the mean of the distribution in the tangent space $T_{\Xm} \P_d$ and $\Sigma$ is the covariance matrix in the tangent space $T_{\Xm} \P_d$.

This distribution has been studied in more depth in \cite{desurrel2025wrappedgaussianmanifoldsymmetric}. It is simple to sample from a wrapped Gaussian $\WG(\Xm, \mu, \Sigma)$, as it solely requires sampling from a Euclidean Gaussian $\N(\mu, \Sigma)$ and applying the exponential map given at \cref{eq:exp_log_riem}. The wrapped Gaussian is illustrated at \cref{fig:wg}.
% \begin{remark}
%     The wrapped Gaussian distribution is not identifiable, as stated in Proposition 4.9 of \cite{desurrel2025wrappedgaussianmanifoldsymmetric}. Indeed, the two parameters $\Xm$ and $\mu$ play an opposing role. In what follows, either $\Xm$ or $\mu$ will be fixed \textit{a priori}, therefore no issues will arise from this behavior.
% \end{remark}
Let us give an important rescaling property of the wrapped Gaussian that will be used in the following:
\begin{proposition}[Rescaling property]
  Let $\Xm \in \P_d, \mu \in \R^D, \sigma > 0$ and $\Zm \sim \WG(\Im_d,0_D, \Im_D)$. Then
  \begin{itemize}
    \item $\Xm^{1/2} \Zm \Xm^{1/2} \sim  \WG(\Xm,0_D, \Im_D),$
    \item $\Exp_{\Im_d}(\Log_{\Im_d}(\Zm) + \mu) \sim \WG(\Im_d,\mu,\Im_D)$
    \item $\Zm^{\sigma} = \exp(\sigma \log \Zm) \sim \WG(\Im_d, 0_D,\sigma^2\Im_D).$
  \end{itemize}
  \label{prop:rescaling_prop}
\end{proposition}
This proposition is a special case of Proposition 4.5 of \citet{desurrel2025wrappedgaussianmanifoldsymmetric}.
Finally, we give a last proposition on the spreading of the wrapped Gaussian when the dimension $d$ rises. The proof is given at \cref{appendix:spread}.
\begin{proposition}
  \label{prop:spread}
  Let $\Zm \sim \WG(\Im_d, 0_D, \Im_D)$, then
  \begin{equation*}
    \begin{aligned}
      \mathbb{E}[\delta(\Im_d, \Zm)] &= \sqrt{2}~\frac{\Gamma\left(\frac{d(d+1)+2}{4}\right)}{\Gamma\left(\frac{d(d+1)}{4}\right)} \\
      &\sim \sqrt{\frac{d(d+1)}{2}} \quad (d \rightarrow +\infty),
    \end{aligned}
  \end{equation*}
  where $\Gamma$ is the Gamma function (see Eq. 16 of \citealt{euler}).
\end{proposition}

\section{Background on Conditional Flow Matching}
\label{sec:background_CFM}
In this section, we recall the classical theory of Conditional Flow Matching (CFM). The general goal of CFM is to sample points from an unknown data distribution $q_1$. In order to do this, CFM learns a parametric map $\Phi$ that pushes forward a ``simple'' distribution $p$ (e.g. a Gaussian) to the data distribution $q_1$. We start by describing the classical Euclidean CFM and then we describe two state-of-the-art approaches to adapt CFM to a Riemannian manifold.

\subsection{The classical Euclidean Conditional Flow Matching}
\label{sec:euc_CFM}
Let us first start by describing the classical Euclidean Conditional Flow Matching (CFM) framework on $\R^d$ from \citet{lipmanFlowMatchingGenerative2023}. Let $p=\mathcal{N}(0,\Im_d)$ be the source distribution and $q_1$ be the target distribution. In order to build a model capable of generating new samples from $q_1$ on $\R^d$, the Flow Matching framework builds a \emph{probability path} $(p_t)_{0 \leq t \leq 1}$ such that $p_0 = p$ and $p_1 = q_1$. To model it, a time dependent flow $\Phi: [0,1] \times \mathbb{R}^d \longrightarrow \mathbb{R}^d$ is defined as the solution of the ODE driven by a vector field $u_t(\cdot)$ at time $t$ and initial condition $\Phi_0(x) = x$:
\begin{equation}
  \label{eq:ODE}
  \frac{\dd}{\dd t}\Phi_t(x) = u_t\!\left(\Phi_t(x)\right), \quad \Phi_0(x) = x.
\end{equation}
The flow $\Phi$ pushes $p$ to $p_t=\Phi_t \# p$. Thus, transforming samples of $p$ with $\Phi_1$ will give a sample of $q_1$.

Moreover, $p_t$ and $u_t$ are linked by the continuity equation:
\begin{equation}
  \frac{\dd}{\dd t} p_t\left(\Phi_t(x)\right) + \operatorname{div}u_t\left(\Phi_t(x)\right) = 0.
  \label{eq:continuity_eq}
\end{equation}

Assuming that we know how to define $(p_t)_t$ (and thus $u_t$), one could regress the vector field $u_t$ with a neural network $v^\theta_t(x)$ directly:
\begin{equation*}
  \mathcal{L}_{\text{FM}}(\theta) = \mathbb{E}_{\substack{t\sim \mathcal{U}(0, 1)\\ x\sim p_t}} \lVert v_t^\theta(x) - u_t(x) \rVert^2.
\end{equation*}
The main challenge lies in constructing such a flow without access to the explicit form of $p_1$, relying solely on available samples. Therefore, \citet{lipmanFlowMatchingGenerative2023} introduces Conditional Flow Matching (CFM) by designing $p_t$ as a mixture of tractable conditional probability paths $p_t(\cdot\mid z)$ that vary with some conditioning variable $z$:
\begin{equation*}
  p_t(x)=\int p_t\left(x \mid z\right) q\left(z\right) \dd z,
\end{equation*}
where $q(z)$ is a distribution over the conditioning variable. Through the continuity \cref{eq:continuity_eq}, the conditional probability path $p_t(\cdot \mid z)$ defines a conditional vector field $u_t(\cdot\mid z)$, such that the marginalization against $z$ of $u_t(\cdot\mid z)$ generates $p_t$:
\begin{equation*}
  u_t(x)=\int u_t\left(x \mid z\right) \frac{p_t\left(x \mid z\right) q\left(z\right)}{p_t(x)} \dd z.
\end{equation*}
Therefore, rather than regressing $u_t(x)$, we regress the conditional vector field $u_t(x\mid z)$:
\begin{equation*}
  \mathcal{L}_{\text{CFM}}(\theta) = \mathbb{E}_{\substack{t\sim \mathcal{U}(0, 1)\\ z\sim q\\ x\sim p_t(\cdot\mid z)}} \big\| v_t^\theta(x) - u_t(x \mid z)\|^2.
\end{equation*}
The power of the CFM framework lies in the fact that $\mathcal{L}_{\text{CFM}}$ and $\mathcal{L}_{\text{FM}}$ are equal up to a constant (see Theorem 2 of \citealt{lipmanFlowMatchingGenerative2023}). Therefore, their gradients are equal and minimizing $\mathcal{L}_{\text{CFM}}$ will also minimize $\mathcal{L}_{\text{FM}}$.

By choosing a conditional variable $x_1 \sim q_1$ from the target distribution (in practice, $x_1$ is a random training point), and Gaussian probability paths $p_t(x\mid x_1)=\mathcal{N}(x \mid \mu_t(x_1), \sigma_t^2 \operatorname{Id})$, the conditional vector field $u_t\left(x \mid x_1\right)$ is given by (see \citealt{lipmanFlowMatchingGenerative2023}):
\begin{equation*}
  u_t\left(x \mid x_1\right)=\frac{\sigma_t^{\prime}\left(x_1\right)}{\sigma_t\left(x_1\right)}\left(x-\mu_t\left(x_1\right)\right)+\mu_t^{\prime}\left(x_1\right).
\end{equation*}
\subsection{Riemannian Conditional Flow Matching}
\label{sec:riem_CFM}
Now that we have explained how Euclidean CFM works, let us describe two possible ways of adapting it when the data distribution $q_1$ lives on the Riemannian manifold $\P_d$ of SPD matrices. We start by describing \emph{RiemannianCFM} introduced in \citet{chen2024flow} that relies on geodesic interpolation and then, we review \emph{DiffeoCFM}, from \citet{collasRiemannianFlowMatching2025} that leverages the tangent spaces.
\paragraph{RiemannianCFM}
In \citet{chen2024flow}, the authors propose to adapt the Euclidean CFM described in \cref{sec:euc_CFM} to the case of a Riemannian manifold.
For this, the conditioning variable $z$ is a pair $(\Xm_0, \Xm_1) \sim q = p \otimes q_1$ sampled from the source distribution $p$ and the target distribution $q_1$ (in practice, $\Xm_1$ is uniformly chosen among the training dataset). Then, they leverage geodesics to build the conditional probability paths $p^{\text{RCFM}}_t(\cdot \mid \Xm_0, \Xm_1)$, as they are the shortest paths between two points on a manifold. The resulting conditional probability path $p^{\text{RCFM}}_t(\cdot \mid \Xm_0, \Xm_1)$ and its associated flow are:
\begin{equation}
  \begin{aligned}
    &p^{\text{RCFM}}_t(\cdot \mid \Xm_0, \Xm_1) = \delta_{\gamma_{\Xm_0,\Xm_1}(t)} \\
    &\Phi^{\text{RCFM}}_t(\Xm\mid \Xm_0, \Xm_1) = \gamma_{\Xm_0,\Xm_1}(t) = \Exp_{\Xm_0}(t\Log_{\Xm_0}\Xm_1),
  \end{aligned}
\end{equation}
where $\delta_{\Xm}$ is the Dirac distribution at $\Xm$ and $\gamma_{\Xm_0,\Xm_1}$ is the unique geodesic from $\Xm_0$ to $\Xm_1$. Then, the conditional vector fields $u^{\text{RCFM}}_t(\cdot \mid \Xm_0,\Xm_1)$ that generate $p^{\text{RCFM}}_t(\cdot \mid \Xm_0, \Xm_1)$ is simply $\frac{\dd}{\dd t} \gamma_{\Xm_0,\Xm_1}(t)$.
\paragraph{DiffeoCFM}
The authors of \citet{collasRiemannianFlowMatching2025} use a Euclidean space $E$ to pullback the distributions $p$ and $q_1$ onto $E$ and then rely on classical Euclidean CFM in E as described in \cref{sec:euc_CFM}. When the space of interest is $\P_d$, the chosen Euclidean space $E$ is the tangent space at identity $T_{\Im_d} \P_d$ and the pullback is the matrix logarithm $\log = \Log_{\Im_d}$. Thus, by denoting $p^{\text{E}}_t$ the Euclidean conditional probability path and by $\Phi^{\text{E}}_t$ the Euclidean flow (as given in \citealt{lipmanFlowMatchingGenerative2023}), we can express them for \emph{DiffeoCFM}:
\begin{equation}\label{eq:diffeoCFM}
  \begin{aligned}
    &p^{\text{DiffeoCFM}}_t(\cdot \mid \Xm_1) = \exp \# p^{\text{E}}_t(\cdot \mid \log\Xm_1) \\
    &\Phi^{\text{DiffeoCFM}}_t(\cdot \mid \Xm_1) = \exp \circ \Phi^{\text{E}}_t(\cdot \mid \log \Xm_1)
  \end{aligned}
\end{equation}
In this case, the conditional variable $z$ is the $\log$ of a point $\Xm_1 \sim q_1$. One can also express the conditional vector field generating $p^{\text{DiffeoCFM}}_t(\cdot \mid \Xm_1)$ (see Equation (7) of \citealt{collasRiemannianFlowMatching2025}):
\begin{equation}\label{eq:diffeo_cfm_u_cond}
  u^{\text{DiffeoCFM}}_t(\Xm \mid \Xm_1) = \dd \exp(\log\Xm)\cdot\left[u^{\text{E}}_t(\log\Xm)\right].
\end{equation}
where $\dd \exp(\Um)\cdot \Vm$ is the differential of $\exp$ at $\Um$ applied at $\Vm$.
\section{A global framework unifying Riemannian FM using the wrapped Gaussian}
\label{sec:global_framework}
Now that we have described the general approach of CFM and two specific Riemannian methods, \emph{RiemannianCFM} and \emph{DiffeoCFM}, we propose here a very general framework for doing CFM on the Riemannian manifold $\P_d$ of SPD matrices. We will see that this framework encompasses \emph{RiemannianCFM} and \emph{DiffeoCFM} but can also propose new Riemannian CFM methods. This framework is partially inspired by the one provided in \citet{tong2024improving}.
\subsection{The general framework}
This framework relies on the wrapped Gaussian distribution as described at \cref{sec:wrapped_gaussian}. In this framework, the source distribution $p$ will be a simple centered wrapped Gaussian: $p = \WG(\Im_d,0_D, \sigma_{\text{init}}^2\Im_D)$. This distribution is convenient because of its ease to sample from (see \cref{sec:wrapped_gaussian}). Let us now consider a conditioning variable $\Zm \sim q$, where $q(\Zm)$ is a distribution over the conditioning variable. We consider the following conditional probability path $\left(p_t(\cdot \mid \Zm)\right)_{0 \leq t \leq 1}$:
\begin{equation}
  p_t(\cdot \mid \Xm_1) = \WG(\bar{\Xm}_t(\Zm), \mu_t(\Zm), \sigma_t(\Zm)^2 \Im_D),
  \label{eq:general_conditional_proba_path}
\end{equation}
where $\bar{\Xm} \colon [0,1] \times \P_d \rightarrow \P_d$, $\mu \colon [0,1] \times \P_d \rightarrow \R^D$ and $\sigma \colon [0,1] \times \P_d \rightarrow \R_{>0}$ parametrize the conditional probability path. This conditional probability path $(p_t(\cdot \mid \Xm_1))_t$ is defined directly on the manifold $\P_d$ and relies on a infinitely supported distribution. In order to match the source distribution $p$ at $t=0$, we impose that the marginal $p_0 = p$. This marginal $p_0$ can be computed as $p_0 =  \int p_0(\cdot \mid \Zm) q(\Zm)\dd\Zm$, which constrains jointly the conditional path $p_0(\cdot \mid \Zm)$ and the distribution $q$ over the conditioning variable $\Zm$.
%$$\bar{\Xm}_0(\Zm) = \Im_d,~ \mu_0(\Zm) = 0_D,~ \text{and }  \sigma_0(\Zm) =\sigma_{\text{init}}.$$
%It would nevertheless be possible to consider another wrapped Gaussian for the source distribution, leading to other initial conditions. We make this choice to simplify the following computations.
\begin{remark}
  Using \cref{prop:spread}, one can see that the spread of $\WG(\Im_d,0_D,\Im_D)$ scales roughly like $d$. This can lead to some numerical issues when $d$ is large. This is why, to solve this issue, we start from $p_0 = \WG(\Im_d, 0_D, \sigma_{\text{init}}^2\Im_D)$ with $\sigma_{\text{init}} <1$, thus a more concentrated source distribution.  In the experiments, we will set $\sigma_{\text{init}} = \sqrt{\frac{2}{d(d+1)}}$.
\end{remark}

We now need to build a conditional flow $\Phi_t$ (conditioned by $\Zm$) such that $\Phi_t \# p_0 = p_t(\cdot \mid \Zm)$. Using the properties of the wrapped Gaussian given in \cref{prop:rescaling_prop}, we have the following proposition (see \cref{appendix:the_maths_general_framework} for the proof):
\begin{proposition}
  The following conditional flow $\Phi_t$ satisfies $\Phi_t \# p_0 = p_t(\cdot \mid \Zm)$ where $p_t(\cdot \mid \Zm)$ is given at \cref{eq:general_conditional_proba_path}
  \begin{equation}\label{eq:general_flow}
    \Phi_t(\Xm) =
    \bar{\Xm}_t^{1/2}\exp\left[\frac{\sigma_t}{\sigma_{\text{init}}}\log(\Xm) + \mu_t\right]\bar{\Xm}_t^{1/2},
  \end{equation}
  where we dropped the dependency in $\Zm$ for all time-dependent quantities.
\end{proposition}
This conditional flow $\Phi_t$ can be decomposed into two parts: the first one is inside the $\exp$, where one retrieves the classical Euclidean flow between two Gaussians (up to the change of variables $\log(\Xm)$) as given at Equation 11 of \citet{lipmanFlowMatchingGenerative2023}, this part affects the parameters of the Gaussian in the tangent spaces. Then, the second part is the outside of the $\exp$, that controls and adapts the base point of the wrapped Gaussian defining the conditional probability path (see the first point of \cref{prop:rescaling_prop}).

This flow then provides a vector field that generates the conditional probability path $(p_t(\cdot \mid \Zm))_t$:
$$\frac{\dd}{\dd t} \Phi_t(\Xm) = u_t(\Phi_t(\Xm) \mid \Zm) \in T_{\Phi_t(\Xm)} \P_d.$$
One can solve this ODE for $u_t$ to get a closed form of the conditional vector field:

\begin{theorem}
  \label{theo:general_u_cond}
  Let $p_t(\cdot \mid \Zm)$ be the conditional probability path given at \cref{eq:general_conditional_proba_path}
  and $\Phi_t$ its flow. The vector field defining $\Phi_t$ depends on the value of $\sigma_t(\Zm)$ in the following manner:
  \begin{itemize}
    \item If $\sigma_t(\Zm) > 0$, then the unique vector field defining $\Phi_t$ is
      \begin{equation*}
        \begin{aligned}
          u_t(\Ym \mid \Zm)
          &= \Sm(t)\bar{\Xm}_t^{-1/2}\Ym + \Ym\bar{\Xm}_t^{-1/2}\Sm(t) \\
          &\quad + \bar{\Xm}_t^{1/2}
          \,\dd\exp(\log \bm M_t)\cdot
          \left[
            \frac{\sigma_t'}{\sigma_t}(\log \bm M_t - \mu_t)
            + \mu_t'
          \right]
          \bar{\Xm}_t^{1/2},
        \end{aligned}
      \end{equation*}
      where $\Sm$ is the derivative of $t \mapsto \bar{\Xm}_t(\Zm)^{1/2}$, $\bm M_t = \bar{\Xm}_t^{-1/2}\Ym\bar{\Xm}_t^{-1/2}$, and all
      time-dependent quantities are evaluated at $\Zm$.
    \item If $\sigma_t(\Zm) = 0$, then
      \[
        u_t(\Ym \mid \Zm) = \frac{\dd}{\dd t}\left[\bar{\Xm}_t^{1/2}\exp(\mu_t)\bar{\Xm}_t^{1/2}\right],
      \]
      where all time-dependent quantities are evaluated at $\Zm$.
  \end{itemize}

\end{theorem}
The proof can be found at \cref{appendix:the_maths_general_framework}. Let us comment on this result. First, the distinction between the cases $\sigma_t(\Zm) = 0$ and $\sigma_t(\Zm) \neq 0$ is due to the fact that when $\sigma_t(\Zm) = 0$, the conditional flow $\Phi_t$ defined in \cref{eq:general_flow} is not invertible. Second, the time derivative $\Sm$ of $t \mapsto \bar{\Xm}_t(\Zm)^{1/2}$ is the solution to the following Sylvester equation:
(see Theorem 1.26 or Chapter 6 of \citealt{highamFunctionsMatrices2008}):
\begin{equation}
  \label{eq:sylvester_eq}
  \bar{\Xm}_t(\Zm)^{1/2} \,\Sm(t) + \Sm(t)\, \bar{\Xm}_t(\Zm)^{1/2} = \frac{\dd}{\dd t}\bar{\Xm}_t(\Zm),
\end{equation}
which can be numerically solved. However, we will see that for the specific choices of $\bar{\Xm}_t(\Zm)$ that we will use in what follows, $\Sm$ can be explicitly computed. Third, the direction of the $\exp$ differential is exactly the conditional vector field for the Euclidean CFM as given at Theorem 3 of \cite{lipmanFlowMatchingGenerative2023}.  Finally, recall that the vector $u_t(\Xm \mid \Zm)$ belongs to the tangent space $T_{\Xm} \P_d = \S_d$. One can verify that $u_t(\Xm \mid \Zm)$ is indeed symmetric.

\subsection{Specific parameterizations}
\begin{table}[t]
  \centering
  \small
  \setlength{\tabcolsep}{6pt}
  \begin{tabular}{lcccc}
    \toprule
    \textbf{Method}
    & $q(\Zm)$
    & $\bar{\Xm}_t(\Xm_1)$
    & $\mu_t(\Xm_1)$
    & $\sigma_t(\Xm_1)$
    \\
    \midrule
    RiemannianCFM
    & $p(\Xm_0)q_1(\Xm_1)$
    & $\gamma_{\Xm_0,\Xm_1}(t)$
    & $0_D$
    & $0$
    \\

    DiffeoCFM
    & $q_1(\Xm_1)$
    & $\Im_d$
    & $t\log(\Xm_1)$
    & $\sigma_{\text{lin}}$
    \\

    WrappedCFM
    & $q_1(\Xm_1)$
    & $\Xm_1^{t}$
    & $0_D$
    & $\sigma_{\text{lin}}$
    \\
    \bottomrule
  \end{tabular}%

  \caption{Unified Wrapped-Gaussian framework for Riemannian CFM on $\P_d$. Each method corresponds to a specific parameterization of the conditional path $p_t(\cdot\mid \Xm_1)=\mathrm{WG}(\bar{\Xm}_t,\mu_t,\sigma_t^2 I_D)$. We denote the linear scheduler $\sigma_{\text{lin}} = \sigma_{\text{init}} - \left(\sigma_{\text{init}} - \sigma_{\min}\right)t$}
  \label{tab:wg_framework}
\end{table}

The newly built general Riemannian CFM framework for the manifold $\P_d$ relies on the parameterization of the conditional probability paths $(p_t(\cdot \mid \Zm))_t$ by the three following mappings: $\bar{\Xm}$, $\mu$ and $\sigma$ and by the choice of the distribution over the conditioning variable $\Zm$. We will  now see that, depending on the choices of these parameterizations, we are able to retrieve the SOTA methods described in \cref{sec:background_CFM}: the \emph{RiemannianCFM} and the \emph{DiffeoCFM}.
\paragraph{RiemannianCFM} Let us start with the RiemannianCFM. In order to retrieve this method from our general framework, we chose the conditioning variable $\Zm$ to be a pair of a random source point $\Xm_0 \sim p$ and a random target point $\Xm_1 \sim q_1$ thus, $q(\Zm) = p(\Xm_0)q_1(\Xm_1)$ is the independent coupling. Then, the choice for the parameterization of the conditional probability paths is:
\[        \bar{\Xm}_t(\Xm_0, \Xm_1) = \gamma_{\Xm_0, \Xm_1}(t),~
  \mu_t(\Xm_0, \Xm_1) = 0_D, \text{ and }
  \sigma_t(\Xm_0, \Xm_1) = 0,
\]

where we recall that $\gamma_{\Xm_0, \Xm_1}$ is the unique geodesic joining $\Xm_0$ and $\Xm_1$. This is the degenerated case where the covariance of the wrapped Gaussian is zero, leading to a Dirac distribution on $\P_d$. It is similar to the Euclidean setting as described in \citet{tong2024improving} (in which it is referred to as ``Basic form of CFM".)
\paragraph{DiffeoCFM} In order to get DiffeoCFM from our general framework, the conditioning variable $\Zm$ is a single random point $\Xm_1$ from the target distribution $q_1$, and one chooses the following parameterization:

\[
  \bar{\Xm}_t(\Xm_1) = \Im_d,~
  \mu_t(\Xm_1) = t\log(\Xm_1), \text{ and }
  \sigma_t(\Xm_1) = \sigma_{\text{init}} - \left(\sigma_{\text{init}} - \sigma_{\min}\right)t.
\]
Here, everything is sent to the tangent space $T_{\Im_d} \P_d$ at identity, so the base point $\bar{\Xm}_t(\Xm_1)$ of the WG does not change and then classical CFM is applied in $T_{\Im_d} \P_d$.

In \cref{appendix:getting_back_to_the_methods}, we show how the general formulas given above recover, for each specific method, the correct expressions of the conditional flow $\Phi_t$ and the conditional vector field $u_t(\cdot \mid \Xm_1)$. We summarize the different parameterizations in \cref{tab:wg_framework}.

\section{Wrapped Conditional Flow Matching}
\label{sec:W-CFM}
After introducing a general framework for CFM on the Riemannian manifold $\P_d$ of SPD matrices and showing how state-of-the-art methods can be recovered through specific parameterizations, we introduce in this section a new parameterization that leads to a novel Riemannian CFM model. We call this method \emph{WrappedCFM}. This WrappedCFM relies on a conditional probability path that is denoted $(p^{\text{WCFM}}_t)_{0 \leq t \leq 1}$. Similarly to DiffeoCFM, the conditioning variable $\Zm$ is simply a random point $\Xm_1$ from the target distribution $q_1$. Then,  $(p^{\text{WCFM}}_t)_{t}$ is build using the following parameterization of the general conditional probability paths given at \cref{eq:general_conditional_proba_path}:

\begin{equation}
  \label{eq:wcfm_param}
  \begin{cases}
    \bar{\Xm}_t(\Xm_1) = \Xm_1^t, \\
    \mu_t(\Xm_1) = 0_D, \\
    \sigma_t(\Xm_1) = (1-t)\,\sigma_{\text{init}} + t\,\sigma_{\min},
  \end{cases}
\end{equation}
where $\sigma_{\text{min}} \ll \sigma_{\text{init}}$. Therefore, the conditional probability path at time $t = 0$ is $p^{\text{WCFM}}_0(\cdot \mid \Xm_1) = \WG(\Im_d,0_D, \sigma_{\text{init}}^2\Im_D) = p$, and at time $t = 1$, is a concentrated WG around $\Xm_1$. This conditional probability path $(p^{\text{WCFM}}_t)_{0 \leq t \leq 1}$ is the equivalent of the classical Gaussian conditional probability paths  from \citet{lipmanFlowMatchingGenerative2023} but defined on $\P_d$ using wrapped Gaussians. Remark that $\bar{\Xm}$ is the unique geodesic between $\Im_d$ and $\Xm_1$ under the AIRM metric given at \cref{eq:AIRM_metric}.

The differences between the proposed WrappedCFM and RiemannianCFM is that the WrappedCFM relies on an infinitely supported distribution, whereas the RiemannianCFM relies on Dirac distributions (recall that RiemannianCFM is a degenerate case of our framework, when $\sigma_t(\Xm_1) = 0$). The difference between the DiffeoCFM and the WrappedCFM is that the DiffeoCFM uses only \emph{one} tangent space to project the data onto. This tangent space $T_{\Im}\P_d$ is a good local approximation of the manifold $\P_d$, but if the target distribution $q_1$ is very spread out, relying only on $T_{\Im}\P_d$ can induce some deformations. In our proposed WrappedCFM, as the center $\bar{\Xm}_t(\Xm_1)$ of the wrapped Gaussian modeling the conditional probability path $p_t(\cdot \mid \Xm_1)$ changes at each time step $t$, it provides a better approximation of the curved space $\P_d$.

We now need to build a conditional flow $\Phi_t$ (conditioned by $\Xm_1$) such that $\Phi_t \# p^{\text{WCFM}}_0 = p^{\text{WCFM}}_t(\cdot \mid \Xm_1)$. Using the general flow given at \cref{eq:general_flow} and the specific parameterization given at \cref{eq:wcfm_param}, we have that the flow $\Phi^{\text{WCFM}}_t$ is:
\begin{equation}
  \Phi^{\text{WCFM}}_t(\Xm) = \Xm_1^{t/2}\Xm^{\frac{\sigma_t(\Xm_1)}{\sigma_{\text{init}}}}\Xm_1^{t/2}.
  \label{eq:flow}
\end{equation}
In this case, $\sigma_t$ is invertible, so one can apply the first case of \cref{theo:general_u_cond} to get the conditional vector field that generates $p^{\text{WCFM}}_t(\cdot \mid \Xm_1)$:
\begin{theorem}
  Let $p^{\text{WCFM}}_t(\cdot \mid \Xm_1)$ be the conditional probability path given of the WrappedCFM and $\Phi^{\text{WCFM}}_t$ its corresponding
  flow map as in \cref{eq:flow}. Then, the unique vector field that defines $\Phi^{\text{WCFM}}_t$ has the form:
  \begin{equation*}
    \begin{split}
      u^{\text{WCFM}}_t(&\Ym \mid \Xm_1)
      = \tfrac{1}{2}\log(\Xm_1)\,\Ym + \tfrac{1}{2}\Ym \log(\Xm_1)\\
      &+ \tfrac{\sigma'_t(\Xm_1)}{\sigma_t(\Xm_1)} \,\Xm_1^{t/2} \bm{M}_t^{1/2}\log(\bm{M}_t)\bm{M}_t^{1/2}\Xm_1^{t/2},
    \end{split}
  \end{equation*}
  where $\bm{M}_t = \Xm_1^{-t/2}\Ym \Xm_1^{-t/2}$. Consequently, $u_t(\cdot \mid \Xm_1)$ generates the wrapped Gaussian path $p_t(\cdot \mid \Xm_1)$.
  \label{eq:conditional_vector_field}
\end{theorem}
The proof is given at \cref{appendix:getting_back_to_the_methods}. As stated after \cref{theo:general_u_cond}, we know explicitly how to compute the derivative $\Sm$ of $t \mapsto \bar{\Xm}_t(\Xm_1) = \Xm_1^t$ using \cref{lemma:diff_exp}. Thus, we do not need to solve numerically the Sylvester equation give at \cref{eq:sylvester_eq} as in the general case.

Finally, the CFM objective is adapted from \citet{chen2024flow} to the Riemannian setting of $\P_d$:
\begin{equation}
  \mathcal{L}_{\text{CFM}}(\theta)
  = \mathbb{E}_{t, \Xm_1, \Xm}
  \|v_t^\theta(\Xm) - u^{\text{WCFM}}_t(\Xm \mid \Xm_1)\|_{\Xm}^2.
\end{equation}
where $t \sim \mathcal{U}(0,1)$, $\Xm_1 \sim q_1$, $\Xm \sim p^{\text{WCFM}}_t(\cdot \mid \Xm_1)$ and $\|\cdot \|_{\Xm}$ is the norm induced by the AIRM given at \cref{eq:AIRM_metric} in the tangent space $T_{\Xm} \P_d$.

Once the vector field $u_t$ is learned using $v_t^\theta$, to generate a sample $\tilde{\Xm}_1 \sim q_1$, one draws an initial sample $\Xm_0 \sim p^{\text{WCFM}}_0 = \WG(\Im_d,0_D, \sigma_{\text{init}}^2\Im_D)$ and integrating the learned vector field $v_t^\theta$ from $t=0$ to $t=1$ using a Riemannian Euler scheme:
\begin{equation}
  \Xm_{k+1} = \Exp_{\Xm_k}(hv_{t_k}^{\theta}(\Xm_k))
  \label{eq:euler_scheme}
\end{equation}
where $h$ is a time step usually set to $\tfrac{1}{100}$ in the experiments and $t_k = kh$. For more insights on solving numerically an ODE on a Riemannian manifold, we refer to \citet{hairer2011solving}.

\section{Synthetic experiments}
\label{sec:synthetic_exp}

\begin{figure}[ht]
  \centering
  \begin{subfigure}[b]{0.49\textwidth}
    \centering
    \includegraphics[width=0.95\textwidth]{figures/swiss_roll_in_cone.pdf}
    %\caption*{Target distribution}
  \end{subfigure}
  \hfill
  \begin{subfigure}[b]{0.49\textwidth}
    \centering
    \includegraphics[width=0.9\textwidth]{figures/figure_8_in_cone.pdf}
    %\caption*{Trajectories of 50 random samples}
  \end{subfigure}
  \caption{Two examples of the result of the proposed WrappedCFM once the flow is learned. The true target distribution is on the top row and the learned one on the bottom row. In these examples (described in \cref{sec:synthetic_exp}), each point represents a $2 \times 2$ SPD matrices, that can thus be represented in 3D. We can see that, qualitatively, the proposed WCFM is able to learn the correct distributions.}
  \label{fig:examples_wgcfm}
\end{figure}

Now that we have introduced a global framework for the different CFM methods using wrapped Gaussians, and proposed a new method for doing CFM on $\P_d$, namely \emph{WrappedCFM}, we want to compare the different methods numerically. In this section we lead synthetic experiments to compare them. Let us start by detailing the contestants, the metrics we will use to evaluate the different methods and the setup of this experiment \footnote{Upon acceptation, our code will be made available on GitHub.}.
\subsection{The contestants}
\label{sec:contestants}
In the following experiments, we will compare our proposed method, WrappedCFM, with the other Riemannian CFM methods that were introduced in \cref{sec:riem_CFM}: \emph{RiemannianCFM} and \emph{DiffeoCFM}. We will also compare these three Riemannian approaches to a classical Euclidean CFM. In this case, the SPD matrices are treated as Euclidean symmetric matrices by vectorizing their upper-triangular part, and the Riemannian geometry is not taken into account.

All of the methods rely on a deep network to learn $v_t^\theta$, the architectures of which are detailed in \cref{appendix:details_synthetic_experiments}.
\subsection{The metrics}
\label{sec:metrics}
Let us now detail the different metrics we will use to evaluate the different models.

\paragraph{Log-likelihood (LL)} The first metric we will use to evaluate the different CFM methods is the log-likelihood of the test set (denoted LL). Indeed, the log-likelihood of a sample $\Xm_1$ can be computed by reversing the ODE \cref{eq:ODE}. We give more details on this at \cref{appendix:log_lik_est}. The higher the likelihood, the better the method. %To evaluate the different methods, we compute the log-likelihood of the testing set and denote it LL.
\paragraph{The $W_2$-distance}
We will also compare the Wasserstein-2 distance between the points sampled and the true test distribution. We recall that the Wasserstein-2 distance between the true sample distribution $q_1$ and the learned one $\tilde{q}_1$ is defined as follows \citep{peyreComputationalOptimalTransport2020}:
$$W_2(q, \tilde{q}) = \left( \inf_{\gamma \in \Pi(q_1, \tilde{q_1})}\mathbb{E}_{(x,y) \sim \gamma}\delta(x,y)^2\right)^{1/2},$$
where $\Pi(q, \tilde{q})$ is the set of couplings of $q_1$ and $\tilde{q}_1$. This $W_2$-distance was computed using POT \citep{flamary2024pot}.

\paragraph{SPDness} We will also compare the number of sampled matrices that are SPD. The three Riemannian methods are build to ensure the SPDness of the sampled matrices, but it is not the case for the Euclidean CFM.

We also measure the $\alpha$-Precision and $\beta$-Recall of the different methods in \cref{appendix:more_results_synthetics}.
\subsection{The datasets}
For this first synthetic experiment, we generated three different synthetic datasets. The first two datasets are composed of $2 \times 2$ matrices and the last one is of higher dimension. The two $2 \times 2$ SPD datasets are inspired by the ones used in \citet{zhoubrainflow}. We call them \emph{Swiss-Roll} and \emph{Figure-8}. %The set $\P_d$ of $2 \times 2$ SPD matrices being of dimension $3$, elements of $\P_d$ can be visualized in a 3D space.
The higher dimensional dataset is composed of $10 \times 10$ SPD matrices, sampled from a random mixture of three wrapped Gaussians. We refer to this last dataset as \textit{Mixture of WG}. Implementation details are given in \cref{appendix:details_synthetic_experiments}.

\subsection{The results}
\input{TMLR/result_expe_synthetics}

We start by giving some qualitative results for the $2 \times 2$ datasets \emph{Swiss-Roll} and \emph{Figure-8} and  by plotting the true and learned distributions at \cref{fig:examples_wgcfm}. As the matrices are $2\times 2$ SPD matrices, they are symmetric and thus of the form $
\begin{pmatrix}a&b\\b&c
\end{pmatrix}$ and one can plot them in 3D by plotting the tuple $(a,b,c)$. We can see that the proposed WrappedCFM method is able to correctly reproduce the true datasets. The learned distributions using our WrappedCFM are noisier than the true distributions, and one can see some outliers. However, the global shape of the two true datasets are well reproduced so it validates our approach.

We also give the quantitative results at table \cref{table:res_exp_synthetic}. One can see that for the two datasets \emph{Swiss-Roll} and  \emph{Figure-8}, our WrappedCFM approach has the highest log-likelihood and the lowest Wasserstein-2 distance. Therefore, our method reproduces the best the true distributions. Moreover, one can see that the Euclidean-CFM struggles to produce $100 \%$ of SPD matrices, which is a big drawback as we want to make sure that all the sampled points belong to the manifold of SPD matrices $\P_d$. For the other methods, as they rely on geometric operations and are built to produce only SPD matrices, we do not encounter any problem. This problem of producing SPD matrices is even more significant in the \emph{WG Mixture} setup, where the EuclideanCFM only produces $50\%$ of $10 \times 10$ SPD matrices. In this setup, DiffeoCFM produces the best results, while our proposed WrappedCFM is slightly behind in terms of $W_2$-distance.

We also report training and sampling times in \cref{appendix:training_times}.

\section{Real-world data experiments}
\label{sec:real_world_experiments}
\subsection{The datasets}
We evaluate our proposed WrappedCFM framework on 3 real-world datasets from different downstream applications: BNCI2014004, BNCI2014002, BNCI2015001 \citep{Leebbcidataset} from Brain Computer Interfaces (BCI), and Textile \citep{MVTec2021} for anomaly detection for textile manufacturing. All these datasets include time series from which we computed covariance matrices. Details on preprocessing are presented in \cref{appendix:datasets_preprocessing} and the datasets are summarized at \cref{tab:datasets_summary}.

\subsection{Evaluation procedure}
To assess the quality of the learned distribution, we use proxy evaluation metrics, namely accuracy and F1-score, within a downstream classification task. A classifier is trained on real data and subsequently applied to generated samples (10000 samples per class) to determine whether the generated data retain sufficient task-specific information for reliable task discrimination. We also compare the Wasserstein-2 distance between generated samples and the true test data. Training details are presented in \cref{appendix:real_training_details}.

\input{ICML/results_datasets}

\subsection{The results}

We report the results in \cref{tab:results_real_datasets}. Our proposed WrappedCFM achieves performance comparable to RiemannianCFM and DiffeoCFM. On BNCI2014004, which is composed of $3 \times 3$ matrices, the Riemannian methods achieve the best results. Our proposed WrappedCFM struggles more in high dimension, as it was the case in the synthetic experiments. Although EuclideanCFM attains the highest numerical scores on these benchmarks, this advantage is likely due to the reduced optimization complexity of learning linear vector fields in Euclidean space, as opposed to the nonlinear high-dimensional constraints imposed by manifold-aware flows. However, it is critical to note that EuclideanCFM lacks the geometric guarantees of our framework. As demonstrated in our synthetic experiments, it fails to ensure that generated samples remain within the SPD cone. In contrast, WrappedCFM inherently respects the manifold structure, offering a mathematically grounded and flexible framework. These findings open compelling research directions toward increasing the expressivity of Riemannian flow models to bridge the remaining gap with the Euclidean baselines.

\section{Conclusion}
%\subsubsection*{References}
%In this paper, we propose a new method to adapt Conditional Flow Matching (CFM) to the Riemannian manifold $\P_d$ of SPD matrices. Our \emph{WrappedCFM} relies on the Wrapped Gaussian, that allows to generalize the Euclidean Gaussian conditional probability paths to conditional probability paths defined on the manifold $\P_d$. Moreover, we provide a general framework for doing CFM on the manifold $\P_d$, using the wrapped Gaussian. We provide general expressions of the conditional vector field that can be applied to different SOTA Riemannian CFM methods. We then tested and validated our method on synthetic and real datasets, comparing it with SOTA methods. Our results suggest that geometry-aware, infinitely supported conditional paths can improve both fidelity and validity  on $\P_d$. A key limitation is the reliance on repeated matrix $\exp$ / $\log$ evaluations, which may be expensive for large $d$. Therefore, addressing scalability and extending the construction to general Hadamard manifolds are promising directions.
In this paper, we introduced \emph{WrappedCFM}, a novel framework for Conditional Flow Matching to the Riemannian manifold $\P_d$ of SPD matrices. By leveraging the Wrapped Gaussian distribution, our approach provides a unified perspective that encompasses existing state-of-the-art Riemannian CFM methods as special cases while bypassing the need for extrinsic approximations or restrictive pullbacks. Leveraging this framework, we introduce a novel interpolation scheme validated through extensive experiments on both synthetic and real-world datasets. Our results demonstrate that geometry-aware, infinitely supported conditional paths is able to sample diverse samples while maintaining high fidelity. A key limitation remains the computational cost of repeated matrix exponential and logarithm evaluations, particularly in high dimensions. Future work will focus on addressing this scalability and extending our Wrapped Gaussian framework to broader classes of Hadamard manifolds \cite{shigaHadamardManifolds1984} as wrapped Gaussian can be defined in a similar manner. We will also explore new parameterizations of our general probability paths, inspired by recent advances in Euclidean flow-based models \citep{albergo2023stochasticinterpolants,liu2022flow,DBLP:conf/icml/PooladianBDALC23,pmlr-v238-tong24a}, to generalize other CFM methods to the Riemannian manifold $\P_d$.
\bibliography{main}
\bibliographystyle{tmlr}

\appendix

\crefalias{section}{appendix}
\onecolumn

\section{More details on Wrapped Gaussian}
In this appendix, we give more results on the wrapped Gaussian $\WG(\Xm, \mu, \Sigma)$. These results come from \citep{desurrel2025wrappedgaussianmanifoldsymmetric}.
\begin{theorem}
  The density $f_{\Xm,\mu,\Sigma}$ of the wrapped Gaussian $\WG(\Xm, \mu, \Sigma)$ exists and is:
  \begin{equation}
    \label{eq:density}
    \begin{aligned}
      &\forall \Zm \in \P_d,~f_{\Xm,\mu,\Sigma}(x) = \frac{g_{\mu,\Sigma}(\Log_{\Xm}(\Zm))}{|J_{\Xm}(\Log_{\Xm}(\Zm))|}
    \end{aligned}
  \end{equation}
  where $g_{\mu, \Sigma}$ is the density of the multivariate Gaussian $\mathcal{N}(\mu, \Sigma)$ and $J_{\Xm}(\cdot ) = \det(\mathrm{d} \Exp_{\Xm}(\cdot))$ is the Jacobian determinant of the exponential map $\Exp_p$.
  \label{theo:density}
\end{theorem}

The Jacobian determinant of the exponential map describes the volume change induced by the identification of $T_{\Xm}\P_d$ and $\P_d$ by the exponential map $\Exp_{\Xm}$. To compute the density explicitly, one needs to compute the Jacobian $J_{\Xm}(u)$, which has a closed form expression for the manifold $\P_d$ (see Appendix C of  \citealt{desurrel2025wrappedgaussianmanifoldsymmetric}).
\begin{proposition}
  \label{prop:jac_exp}
  The Jacobian determinant of the exponential map at the identity $\Exp_{\Im_d}$ is:
  $$\forall \Um \in T_{\Im_d} \P_d,~ J_{\Im_d}(\Um) = 2^{\nicefrac{d(d-1)}{2}} \prod_{i < j} \frac{\sinh\left(\frac{\lambda_i(\Um) - \lambda_j(\Um)}{2}\right)}{\lambda_i(\Um) - \lambda_j(\Um)}$$
  where the $\lambda_i(\Um)$ are the eigenvalues of $u$.
  Then, one can use the previous formula to compute the Jacobian determinant of the exponential map at any point $\Xm \in \P_d$:
  $$\forall \Um \in T_{\Xm} \P_d,~ J_{\Xm}(\Um) = J_{I_d}(\Xm^{-1/2}\Um\Xm^{-1/2}).$$
\end{proposition}

Let us give a proof of \cref{prop:spread} that we recall here:
\begin{proposition}
  Let $\Zm \sim \WG(\Im_d, 0_D, \Im_D)$, then
  \begin{equation*}
    \begin{aligned}
      \mathbb{E}[\delta(\Im_d, \Zm)] &= \sqrt{2}~\frac{\Gamma\left(\frac{d(d+1)+2}{4}\right)}{\Gamma\left(\frac{d(d+1)}{4}\right)} \\
      &\sim \sqrt{\frac{d(d+1)}{2}} \quad (d \rightarrow +\infty),
    \end{aligned}
  \end{equation*}
  where $\Gamma$ is the Gamma function (see Eq. 16 of \citealt{euler}).
\end{proposition}

In order to prove this, let us start by the Euclidean case.
\begin{proposition}
  \label{prop:spread_gaussian}
  Let $Z \sim \mathcal{N}(0_d, \Im_d)$, then
  \begin{equation*}
    \begin{aligned}
      \mathbb{E}[\|Z\|_2] &= \sqrt{2}~\frac{\Gamma\left(\frac{d+1}{2}\right)}{\Gamma\left(\frac{d}{2}\right)} \\
      &\sim \sqrt{d} \quad (d \rightarrow +\infty),
    \end{aligned}
  \end{equation*}
\end{proposition}
\begin{proof}
  As $Z \sim \mathcal{N}(0_d, \Im_d)$, its norm $\|Z\|_2$ follows a Chi distribution. According to Section 11.3 of \citet{Evans2001StatisticalDT}, its expectation is
  $$\mathbb{E}[\|Z\|_2] =  \sqrt{2}~\frac{\Gamma\left(\frac{d+1}{2}\right)}{\Gamma\left(\frac{d}{2}\right)}.$$
  Moreover, using Stirling's formula, one gets, according to formula 21 of \cite{euler} that
  $$\Gamma(x) \sim x^{x-1/2}e^{-x}\sqrt{2\pi}.$$
  Thus, denoting $n = \frac{d}{2}$, on gets
  $$\mathbb{E}[\|Z\|_2] \sim \sqrt{2}\frac{(n+\frac{1}{2})^{n}e^{-(n+\frac{1}{2})}\sqrt{2\pi}}{n^{n - \frac{1}{2}} e^{-n}\sqrt{2\pi}} = \sqrt{2}e^{-1/2} \sqrt{n}(1+\frac{1}{2n})^n \sim \sqrt{2}\sqrt{n} = \sqrt{d}.$$
\end{proof}
Going back to the wrapped Gaussian case:
\begin{proof}[Proof of Proposition 3.4]
  Let $\Zm \sim \WG(\Im_d, 0_D, \Im_D)$.
  Then,
  \begin{equation*}
    \begin{aligned}
      \mathbb{E}[\delta(\Im_d, \Zm)] &= \int_{\P_d} \delta(\Im_d, \Zm) f_{(\Im_d, 0_D, \Im_D)}(\Zm)\dd\text{vol}(\Zm) \\
      &= \int_{\P_d} \|\Log_{\Im_d}(\Zm)\|_{\Im_d} ~f_{(\Im_d, 0_D, \Im_D)}(\Zm)\dd\text{vol}(\Zm) \\
      &= \int_{\R^{\frac{d(d+1)}{2}}} \|t\|_2 ~g_{0_D, \Im_D}(t) \dd t \quad (\text{where } t = \Log_{\Im_d}(\Zm))\\
      &= \mathbb{E}[\|t\|_2]
    \end{aligned}
  \end{equation*}
  where $f_{(\Im_d, 0_D, \Im_D)}$ is the density of the wrapped Gaussian $\WG(\Im_d, 0_D, \Im_D)$ given at \cref{theo:density}, and $t \sim \mathcal{N}(0_D, \Im_D)$ and $g_{0_D, \Im_D}$ is the density of $\mathcal{N}(0_D, \Im_D)$. Thus, we can use the \cref{prop:spread_gaussian} to conclude:
  $$\mathbb{E}[\delta(\Im_d, \Zm)] \sim \sqrt{D} = \sqrt{\frac{d(d+1)}{2}}.$$
\end{proof}
\label{appendix:spread}
% \section{Mathematical proofs for the Wrapped Conditional Flow Matching}
% \label{appendix:the_maths}
% In the section, we derive the mathematical proofs of the different statements given in \cref{sec:W-CFM}.

% \paragraph{} At \cref{eq:flow}, we give the following conditional flow $\Phi_t$:
% $$\Phi_t(\Xm) = \Xm_1^{t/2}\Xm^{\frac{\sigma_t(\Xm_1)}{\sigma_{\text{init}}}}\Xm_1^{t/2},$$
% and we state the this flow generates the conditional probability path given at \cref{eq:conditional_proba_path}. Indeed, using the second point of \cref{prop:rescaling_prop} and given $\Xm \sim p = \WG(\Im_d,0_D, \sigma_{\text{init}}^2\Im_D), \Xm_1 \sim q$ and $t \in [0,1]$, we have that
% $$\Xm^{\frac{\sigma_t(\Xm_1)}{\sigma_{\text{init}}}} \sim \WG\left(\Im_d, 0_D,\sigma_{\text{init}}^2\frac{\sigma_t(\Xm_1)^2}{\sigma_{\text{init}}^2} \Im_D\right) = \WG(\Im_d, 0_D,\sigma_t(\Xm_1)^2 \Im_D).$$
% Then, using the first point of \cref{prop:rescaling_prop}, we get that
% $$\Xm_1^{t/2}\Xm^{\sigma_t(\Xm_1)}\Xm_1^{t/2} \sim \WG(\Xm_1^{t},0_D \sigma_t(\Xm_1)^2 \Im_D),$$
% which is exactly $p_t(\cdot \mid \Xm_1)$. Thus, we have that $\Phi_t \# p = p_t(\cdot \mid \Xm_1)$.

% \paragraph{} Let us now prove \cref{eq:conditional_vector_field}. We start by recalling it:
% \begin{theorem}
%     Let $p(\cdot \mid \Xm_1)$ be the conditional probability path given at \cref{eq:conditional_proba_path} and $\Phi_t$ its corresponding
% flow map as in \cref{eq:flow}. Then, the unique vector field that defines $\Phi_t$ has the form:
% \begin{equation}
%     u_t(\Ym \mid \Xm_1)
%   = \tfrac{1}{2}\log(\Xm_1)\,\Ym
%   + \tfrac{\sigma'_t(\Xm_1)}{\sigma_t(\Xm_1)} \, \Xm_1^{t/2} \bm{M}^{1/2}\log(\bm{M})\bm{M}^{1/2}\Xm_1^{t/2}
%   + \tfrac{1}{2}\Ym \log(\Xm_1),
% \end{equation}
% where $\bm{M} = \Xm_1^{-t/2}\Ym \Xm_1^{-t/2}$. Consequently, $u_t(\cdot \mid \Xm_1)$ generates the Gaussian path $p_t(\cdot \mid \Xm_1)$.
% \end{theorem}
% In what follows, we set $\sigma_{\text{init}} = 1$ to simplify the computations. We will comment on the general result after the proof. Before proving this theorem, let us give two lemmas:
% \begin{lemma}[Equation 10.3 of \citet{serre2002matrices}]
%     Let $A \in \R^{d \times d}$, then
%     $$\frac{\d}{\d t} \exp(tA) = A\exp(tA) = \exp(tA)A.$$
%     \label[lemma]{lemma:diff_exp}
% \end{lemma}
% % \begin{lemma}
% %     Let $t \in [0,1]$ and $\Xm,\Xm_1 \in \P_d$, then,
% %     $$\Xm^{-\frac{1}{2}\sigma_t(\Xm_1)}\log(\Xm) = \log(\Xm)\Xm^{-\frac{1}{2}\sigma_t(\Xm_1)}.$$
% %     \label{lemma:commut}
% % \end{lemma}
% % \begin{proof}[Proof of \cref{lemma:commut}]
% %     Using the previous lemma, one has:
% %     $$\frac{\d}{\d t} \Xm^{-\frac{1}{2}\sigma_t(\Xm_1)} = \frac{\d}{\d t} \exp\left(-\frac{1}{2}\sigma_t(\Xm_1) \log(\Xm)\right) = -\frac{1}{2}\sigma'_t(\Xm_1)\log(\Xm) \Xm^{-\frac{1}{2}\sigma_t(\Xm_1)} = -\frac{1}{2}\sigma'_t(\Xm_1)\Xm^{-\frac{1}{2}\sigma_t(\Xm_1)}\log(\Xm).$$
% %     Thus, it gives us (as $\sigma'_t(\Xm_1) \neq 0$)
% %     $$\Xm^{-\frac{1}{2}\sigma_t(\Xm_1)}\log(\Xm) = \log(\Xm)\Xm^{-\frac{1}{2}\sigma_t(\Xm_1)}.$$
% % \end{proof}

% \begin{lemma}
%     \label[lemma]{lemma:commut}
%     Let $\Zm \in \P_d$ and $\alpha \in \R$, then $\Zm^\alpha\log(\Zm) = \log(\Zm)\Zm^\alpha$.
% \end{lemma}
% \begin{proof}[Proof of \cref{lemma:commut}]
%     Let us consider the spectral decomposition of $\Zm$: $\Zm = \Um \Lambda \Um^\top,$ where $\Um$ is an orthogonal matrix and $\Lambda$ is a diagonal matrix. Then, on has that $\log \Zm = \Um \log(\Lambda) \Um^\top$ and $ \Zm^\alpha = \Um \Lambda^\alpha \Um^\top$. Thus
%     $$\Zm^\alpha\log(\Zm) = \Um \Lambda^\alpha \Um^\top \Um \log (\Lambda) \Um^\top =   \Um \Lambda^\alpha \log (\Lambda) \Um^\top =  \Um \log (\Lambda) \Lambda^\alpha \Um^\top = \log(\Zm)\Zm^\alpha.$$
%     The third equality coming from the fact that $\Lambda^\alpha$ and $\log(\Lambda)$ are diagonal so they commute.
% \end{proof}
% Let us now go back to the proof of \cref{eq:conditional_vector_field}.
% \begin{proof}[Proof of \cref{eq:conditional_vector_field}]
%     In order to prove the expression of this vector field, we need to go back to the following ODE:
%     $$\frac{\d}{\d t} \Phi_t(\Xm) = u_t(\Phi_t(\Xm) \mid \Xm_1).$$
%     We thus need to compute $\frac{\d}{\d t} \Phi_t(\Xm)$ where $\Phi_t(\Xm)$ is the $\Phi_t$ its conditional flow map given in \cref{eq:flow}. Using \cref{lemma:diff_exp}, one can differentiate each term of the conditional flow $\Phi_t$ as follows:
%     \begin{itemize}
%     \item First,
%     \begin{equation}\label{eq:derivee_x_1_2}
%         \frac{\d}{\d t} \Xm_1^{t/2} = \frac{\d}{\d t} \exp\left(\frac{t}{2}\log(\Xm_1)\right) = \frac{1}{2}\log(\Xm_1)\Xm_1^{t/2} = \frac{1}{2}\Xm_1^{t/2}\log(\Xm_1).
%     \end{equation}
%     \item Similarly, one has
%     $$\frac{\d}{\d t}\Xm^{\sigma_t(\Xm_1)} = \frac{\d}{\d t} \exp(\sigma_t(\Xm_1)\log(\Xm)) = \sigma'_t(\Xm_1)\log(\Xm) \exp(\sigma_t(\Xm_1)\log(\Xm)) =   \sigma'_t(\Xm_1)\log(\Xm)\Xm^{\sigma_t(\Xm_1)}.$$
%     \end{itemize}
%     Thus, on gets
%     \begin{equation*}
%         \begin{aligned}
%             \frac{\d}{\d t} \Phi_t(\Xm) &= \frac{1}{2}\log(\Xm_1)\Phi_t(\Xm) + \sigma'_t(\Xm_1)\Xm_1^{t/2}\log(\Xm)\Xm^{\sigma_t(\Xm_1)}\Xm_1^{t/2} + \frac{1}{2}\Phi_t(\Xm)\log(\Xm_1) \\
%             &= \frac{1}{2}\log(\Xm_1)\Phi_t(\Xm) + \sigma'_t(\Xm_1)\Xm_1^{t/2}\Xm^{\tfrac{1}{2}\sigma_t(\Xm_1)}\log(\Xm)\Xm^{\tfrac{1}{2}\sigma_t(\Xm_1)}\Xm_1^{t/2} + \frac{1}{2}\Phi_t(\Xm)\log(\Xm_1), \\
%         \end{aligned}
%     \end{equation*}
%     where we used \cref{lemma:commut} to get the second equality. Recall that $\Phi_t$ solves the ODE $\frac{\d}{\d t} \Phi_t(\Xm) = u_t(\Phi_t(\Xm) \mid \Xm_1)$.

%     % Let us focus on the middle term:
%     % \begin{equation*}
%     %     \begin{aligned}
%     %         \sigma'_t(\Xm_1)\Xm_1^{t/2}\log(\Xm)\Xm^{\sigma_t(\Xm_1)}\Xm_1^{t/2} &= \sigma'_t(\Xm_1)\Phi_t(\Xm)\Xm_1^{-t/2}\Xm^{-\sigma_t(\Xm_1)}\log(\Xm)\Xm^{-t/2}\Phi_t(\Xm) \\
%     %         &= \sigma'_t(\Xm_1)\Phi_t(\Xm)\Xm_1^{-t/2}\Xm^{-\frac{1}{2}\sigma_t(\Xm_1)}\log(\Xm)\Xm^{-\frac{1}{2}\sigma_t(\Xm_1)}\Xm^{-t/2}\Phi_t(\Xm),
%     %     \end{aligned}
%     % \end{equation*}
%     % where the second equality has been obtained using \cref{lemma:commut}. Then putting things back together, we get:
%     % \begin{equation*}
%     %     \begin{aligned}
%     %         \frac{\d}{\d t} \Phi_t(\Xm) &= \frac{1}{2}\log(\Xm_1)\Phi_t(\Xm) + \sigma'_t(\Xm_1)\Phi_t(\Xm)\Xm_1^{-t/2}\Xm^{-\frac{1}{2}\sigma_t(\Xm_1)}\log(\Xm)\Xm^{-\frac{1}{2}\sigma_t(\Xm_1)}\Xm^{-t/2}\Phi_t(\Xm) + \frac{1}{2}\Phi_t(\Xm)\log(\Xm_1) \\
%     %         &= u_t(\Phi_t(\Xm) \mid \Xm_1).
%     %     \end{aligned}
%     % \end{equation*}
%     Now, letting $\Ym = \Phi_t(\Xm)$, we get from \cref{eq:flow} that $\Xm = \exp\left(\frac{1}{\sigma_t(\Xm_1)}\log(\Xm_1^{-t/2}\Ym\Xm_1^{-t/2})\right)$. Then, plugging this in the previous equation, we get that
%     $$u_t(\Ym \mid \Xm_1)
%   = \tfrac{1}{2}\log(\Xm_1)\,\Ym
%   + \tfrac{\sigma'_t(\Xm_1)}{\sigma_t(\Xm_1)} \, \Xm_1^{t/2} \bm{M}^{1/2}\log(\bm{M})\bm{M}^{1/2}\Xm_1^{t/2}
%   + \tfrac{1}{2}\Ym \log(\Xm_1),$$
%   where $\bm{M} = \Xm_1^{-t/2}\Ym \Xm_1^{-t/2}$. Thus, $u_t(\cdot \mid \Xm_1)$ generates the wrapped Gaussian path $p_t(\cdot \mid \Xm_1)$.
% \end{proof}
% This proves \cref{eq:conditional_vector_field} in the case of $\sigma_{\text{init}} = 1$. Now, set $\tilde{\sigma}_t(\Xm_1) = \frac{\sigma_t(\Xm_1)}{\sigma_{\text{init}}}$. The proof above still holds, with the difference that, in front of the second term, one gets a $\tfrac{\tilde{\sigma}'_t(\Xm_1)}{\tilde{\sigma}_t(\Xm_1)}$. One can easily compute that $$\frac{\tilde{\sigma}'_t(\Xm_1)}{\tilde{\sigma}_t(\Xm_1)} = \frac{\sigma'_t(\Xm_1)}{\sigma_t(\Xm_1)},$$
% and thus, the \cref{eq:conditional_vector_field} remains true when $\sigma_{\text{init}} \neq 1$.

\section{Mathematical details on the general framework}
\label{appendix:the_maths_general_framework}

\subsection{The proof of the general conditional vector field}
In this section, we give a proof of \cref{theo:general_u_cond}, which we start by recalling.
\begin{theorem}
  Let $p_t(\cdot \mid \Xm_1)$ be the conditional probability path  given at \cref{eq:general_conditional_proba_path}
  and $\Phi_t$ its corresponding flow as in \cref{eq:general_flow}.
  The unique vector field defining $\Phi_t$ is
  \begin{equation*}
    u_t(\Ym \mid \Xm_1) = \Sm(t)\bar{\Xm}_t^{-1/2}\Ym+ \Ym\bar{\Xm}_t^{-1/2}\Sm(t)+ \bar{\Xm}_t^{1/2}\,\dd\exp(\log \bm M_t)\cdot\left[\frac{\sigma_t'}{\sigma_t}(\log \bm M_t-\mu_t)+ \mu_t'\right]\bar{\Xm}_t^{1/2},
  \end{equation*}
  where $\Sm$ is the derivative of $t \mapsto \bar{\Xm}_t(\Xm_1)^{1/2}$, $\bm M_t=\bar{\Xm}_t^{-1/2}\Ym\bar{\Xm}_t^{-1/2}$ and all
  time-dependent quantities are evaluated at $\Xm_1$.
\end{theorem}
% Before proving this theorem, let us give two lemmas:
% \begin{lemma}[Equation 10.3 of \citet{serre2002matrices}]
%     Let $A \in \R^{d \times d}$, then
%     $$\frac{\d}{\d t} \exp(tA) = A\exp(tA) = \exp(tA)A.$$
%     \label[lemma]{lemma:diff_exp}
% \end{lemma}
% \begin{lemma}
%     \label[lemma]{lemma:commut}
%     Let $\Zm \in \P_d$ and $\alpha \in \R$, then $\Zm^\alpha\log(\Zm) = \log(\Zm)\Zm^\alpha$.
% \end{lemma}
% \begin{proof}[Proof of \cref{lemma:commut}]
%     Let us consider the spectral decomposition of $\Zm$: $\Zm = \Um \Lambda \Um^\top,$ where $\Um$ is an orthogonal matrix and $\Lambda$ is a diagonal matrix. Then, on has that $\log \Zm = \Um \log(\Lambda) \Um^\top$ and $ \Zm^\alpha = \Um \Lambda^\alpha \Um^\top$. Thus
%     $$\Zm^\alpha\log(\Zm) = \Um \Lambda^\alpha \Um^\top \Um \log (\Lambda) \Um^\top =   \Um \Lambda^\alpha \log (\Lambda) \Um^\top =  \Um \log (\Lambda) \Lambda^\alpha \Um^\top = \log(\Zm)\Zm^\alpha.$$
%     The third equality coming from the fact that $\Lambda^\alpha$ and $\log(\Lambda)$ are diagonal so they commute.
% \end{proof}
% Let us now go back to the proof of \cref{eq:conditional_vector_field}.
\begin{proof}
  We begin by recalling the conditional flow $\Phi_t$ from \cref{eq:general_flow}:
  \begin{equation*}
    \Phi_t(\Xm) = \bar{\Xm}_t(\Xm_1)^{1/2}\exp\left[\frac{\sigma_t(\Xm_1)}{\sigma_{\text{init}}}\log(\Xm) + \mu_t(\Xm_1)\right]\bar{\Xm}_t(\Xm_1)^{1/2}.
  \end{equation*}
  Let us differentiate this flow with respect to $t$. For this, we denote
  \begin{equation*}
    \begin{cases}
      \Sm \text{ the derivative of } t \mapsto \bar{\Xm}_t(\Xm_1)^{1/2} \\
      g(t) = \frac{\sigma_t(\Xm_1)}{\sigma_{\text{init}}}\log(\Xm) + \mu_t(\Xm_1).
    \end{cases}
  \end{equation*}

  Then, one has
  \begin{equation*}
    \frac{\dd}{\dd t} \Phi_t(\Xm) = \Sm(t)\exp\left[g(t)\right]\bar{\Xm}_t(\Xm_1)^{1/2} + \bar{\Xm}_t(\Xm_1)^{1/2}\dd \exp(g(t))\cdot[g'(t)]\bar{\Xm}_t(\Xm_1)^{1/2} +  \bar{\Xm}_t(\Xm_1)^{1/2}\exp\left[g(t)\right]\Sm(t).
  \end{equation*}
  Then, one can compute that
  $$g'(t) = \frac{\sigma'_t(\Xm_1)}{\sigma_{\text{init}}}\log(\Xm) + \mu'_t(\Xm_1).$$
  Thus, doing back to $\frac{\dd}{\dd t} \Phi_t(\Xm)$:
  \begin{align*}
    \frac{\dd}{\dd t} \Phi_t(\Xm) &= \Sm(t)\bar{\Xm}_t^{-1/2} \Phi_t(\Xm) + \Phi_t(\Xm)\bar{\Xm}_t^{-1/2}\Sm(t) \,+\\
    & \bar{\Xm}_t(\Xm_1)^{1/2}\dd \exp\left(\frac{\sigma_t(\Xm_1)}{\sigma_{\text{init}}}\log(\Xm) + \mu_t(\Xm_1)\right)\cdot\left[\frac{\sigma'_t(\Xm_1)}{\sigma_{\text{init}}}\log(\Xm) + \mu'_t(\Xm_1)\right]\bar{\Xm}_t(\Xm_1)^{1/2} \\
    &= u_t(\Phi_t(\Xm) \mid \Xm_1).
  \end{align*}
  Let us denote $\Ym = \Phi_t(\Xm)$. We now need to distinguish between the values of $\sigma_t$.
  \begin{itemize}
    \item \textbf{If} $\sigma_t \neq 0$, then $\Phi_t$ is invertible and $\Xm = \exp\left(\frac{\sigma_{\text{init}}}{\sigma_t(\Xm_1)}\left(\log\bm M_t - \mu_t(\Xm_1)\right)\right)$ where $\bm M_t=\bar{\Xm}_t^{-1/2}\Ym\bar{\Xm}_t^{-1/2}$. Plugging this into the previous equation, we get:
      $$u_t(\Ym \mid \Xm_1) = \Sm(t)\bar{\Xm}_t^{-1/2}\Ym+ \Ym\bar{\Xm}_t^{-1/2}\Sm(t)+ \bar{\Xm}_t^{1/2}\,\dd\exp(\log \bm M_t)\cdot\left[\frac{\sigma_t'}{\sigma_t}(\log \bm M_t-\mu_t)+ \mu_t'\right]\bar{\Xm}_t^{1/2}.$$
    \item \textbf{If} $\sigma_t = 0$, then $\Phi_t$ is not invertible. Going back to the general form of $\Phi_t$ given at \cref{eq:general_flow}, we have that
      $$\Phi_t(\Xm)= \bar{\Xm}_t(\Xm_1).$$
      Thus, the ODE is now simply $$\frac{\dd}{\dd t} \Phi_t(\Xm) = \frac{\dd}{\dd t}\bar{\Xm}_t(\Xm_1) = u_t(\Phi_t(\Xm) \mid \Xm_1),$$
      and thus $u_t(\Phi_t(\Xm) \mid \Xm_1) = \frac{\dd}{\dd t}\bar{\Xm}_t(\Xm_1).$
  \end{itemize}
  Thus, $u_t(\cdot \mid \Xm_1)$ generates the wrapped Gaussian path $p_t(\cdot \mid \Xm_1)$.
\end{proof}

\subsection{Getting back the specific methods}
\label{appendix:getting_back_to_the_methods}
The goal of this section is to show that the general methods, once instantiated with the specific formulas for $\bar{\Xm}_t$, $\mu_t$ and $\sigma_t$, correctly lead to the formulas for the different methods.
\paragraph{DiffeoCFM}
As detailed in \cref{tab:wg_framework}, to instantiate the \emph{DiffeoCFM}, one uses:
$$\Xm_t(\Xm_1) = \Im_d, ~ \mu_t(\Xm_1) = t\log(\Xm_1), \text{ and } \sigma_t(\Xm_1) = \sigma_{\text{init}} - \left(\sigma_{\text{init}} - \sigma_{\min}\right)t.$$
Going back to the conditional flow $\Phi_t$ as given at \cref{eq:general_flow}, we have:
$$\Phi_t(\Xm) = \exp\left[\frac{\sigma_t}{\sigma_{\text{init}}}\log(\Xm) + t \log(\Xm_1)\right],$$
which is exactly the flow as given at \cref{eq:diffeoCFM} with the Euclidean flow being the one from Equation 11 of \citet{lipmanFlowMatchingGenerative2023}. Let us now focus on the conditional vector field given at \cref{theo:general_u_cond}. As $\Xm_t(\Xm_1) = \Im_d$, the time derivative of $\Xm_t(\Xm_1)^{1/2}$ is $0$, thus, $\Sm(t) =0$. Then, the conditional vector field is:
\begin{align*}
    u_t(\Ym \mid \Xm_1) &= \dd \exp(\log(\Ym))\cdot \left[\frac{\sigma_t'}{\sigma_t}(\log \Ym - t \log \Xm_1) + \log \Xm_1 \right] \\ 
                        &= \dd \exp(\log(\Ym))\cdot \left[\frac{\sigma_{\text{init}}\log \Xm_1 - \left(\sigma_{\text{init}} - \sigma_{\min}\right) \log \Ym}{ \sigma_{\text{init}} - \left(\sigma_{\text{init}} - \sigma_{\min}\right)t} \right],
\end{align*}

which is exactly the conditional vector field $u_t^{\text{DiffeoCFM}}$ from \cref{eq:diffeo_cfm_u_cond} where the Euclidean conditional vector field is the one from Equation 21 of \citet{lipmanFlowMatchingGenerative2023}.

\paragraph{WrappedCFM}
Let us do the same work but for \emph{WrappedCFM}, where one uses:
$$\Xm_t(\Xm_1) = \Xm_1^t, ~ \mu_t(\Xm_1) = 0, \text{ and } \sigma_t(\Xm_1) = \sigma_{\text{init}} - \left(\sigma_{\text{init}} - \sigma_{\min}\right)t.$$
Then, the conditional flow is:
$$\Phi_t(\Xm) = \Xm_1^{t/2} \exp\left(\frac{\sigma_t(\Xm_1)}{\sigma_{\text{init}}} \log(\Xm)\right)\Xm_1^{t/2} = \Xm_1^{t/2} \Xm^{\frac{\sigma_t(\Xm_1)}{\sigma_{\text{init}}}}\Xm_1^{t/2},$$
which matches with \cref{eq:flow}. Now for the conditional vector field, we first need the following lemma:
\begin{lemma}[Equation 10.3 of \citet{serre2002matrices}]
  Let $A \in \R^{d \times d}$, then
  $$\frac{\dd}{\dd t} \exp(tA) = A\exp(tA) = \exp(tA)A.$$
  \label[lemma]{lemma:diff_exp}
\end{lemma}
Using this lemma, on gets that
$$\Sm(t) = \frac{\dd}{\dd t}\Xm_1^{t/2} =\frac{\dd}{\dd t} \exp\left(\frac{t}{2} \log(\Xm_1)\right) = \frac{1}{2}\log(\Xm_1)\Xm_1 = \frac{1}{2}\Xm_1 \log(\Xm_1).$$
Moreover, let us consider $\Zm \in \P_d$ and $\Um \in \S_d$ such that they commute. Then, one has that $\dd \exp(\Zm)\cdot [\Um] = \exp(\Zm)\Um$ (see Equation 10.15 of \citealt{highamFunctionsMatrices2008}). Using this, we get:
\begin{equation*}
  \begin{aligned}
    u_t(\Ym \mid \Xm_1) &= \frac{1}{2}\log(\Xm_1)\Xm_1^{1/2}\Xm_1^{-1/2}\Ym + \frac{1}{2}\Ym \Xm_1^{-1/2}\Xm_1^{1/2}\log(\Xm_1) + \Xm_1^{1/2}\dd \exp(\log(\bm M_t))\cdot \left[\frac{\sigma_t'}{\sigma_t}\log\bm M_t\right]\Xm_1^{1/2} \\
    &= \frac{1}{2}\log(\Xm_1)\Ym + \frac{1}{2}\Ym \log(\Xm_1) + \Xm_1^{1/2} \bm M_t \frac{\sigma_t'}{\sigma_t}\log(\bm M_t)\Xm_1^{1/2} \\
    &= \frac{1}{2}\log(\Xm_1)\Ym + \frac{1}{2}\Ym \log(\Xm_1) + \Xm_1^{1/2} \bm M_t^{1/2} \frac{\sigma_t'}{\sigma_t}\log(\bm M_t)\bm M_t^{1/2}\Xm_1^{1/2},
  \end{aligned}
\end{equation*}
which is exactly the conditional vector field from \cref{eq:conditional_vector_field}.

\paragraph{RiemannianCFM}
The case of RiemannianCFM is a bit different, as it is a degenerate case of our framework: $\sigma_t(\Xm_1) = 0$ so the covariance of the wrapped Gaussian is $0_D$. Thus, the conditional flow $\Phi_t$ given at \cref{eq:general_flow} does not depend on the point $\Xm$ anymore:
$$\Phi_t(\Xm) = \gamma_{\Xm_0, \Xm_1}(t)^{1/2}\gamma_{\Xm_0, \Xm_1}(t)^{1/2} = \gamma_{\Xm_0, \Xm_1}(t),$$
where we recall that $\gamma_{\Xm_0, \Xm_1}(t)$ is the unique geodesic joining $\Xm_0$ and $\Xm_1$. Thus, the formula for the conditional vector field from \cref{eq:general_flow} cannot be used directly. However, using the previous formula, one can easily derive
$$u_t(\Ym \mid \Xm_0, \Xm_1) = \frac{\dd}{\dd t} \gamma_{\Xm_0, \Xm_1}(t),$$
which is the conditional vector field of the RiemannianCFM as given in \citet{chen2024flow}.
\section{Log-likelihood computation}
\label{appendix:log_lik_est}
An interesting and important feature of CFM is that one is able to compute the log-likelihood of a generated sample $\Xm_1$. The core idea is to use the continuity equation \ref{eq:continuity_eq} that links the probability path $(p_t)_t$ with the Riemannian divergence of the vector field $u_t$:
$$\frac{\dd}{\dd t} \log(p_t(\Phi_t(\Xm))) + \dive(u_t(\Phi_t(\Xm))) = 0.$$
Integrating this continuity equation for $t \in [0,1]$, one can show that the log probability can be computed together with the flow trajectory by solving the ODE:
$$\frac{\dd}{\dd t}
\begin{bmatrix}
  f(t) \\ g(t)
\end{bmatrix} =
\begin{bmatrix}
  u_t(f(t)) \\ - \dive(u_t(f(t)))
\end{bmatrix} \text{ with }
\begin{bmatrix}
  f(1) \\ g(1)
\end{bmatrix} =
\begin{bmatrix}
  \Xm_1 \\ 0
\end{bmatrix}$$
By solving the ODE backwards (from $t = 1$ to $t = 0$), one can compute the following log-likelihood:
$$\log p_1(\Xm_1) = \log p_0 (f(0)) - g(0)$$
where $p_0$ is the source distribution. This is done using a Riemannian Euler scheme similar to the on given at \cref{eq:euler_scheme} with 1000 steps. As our proposed WrappedCFM is instantiated on the Riemannian manifold $\P_d$, we are able to compute the Riemannian divergence of a vector field $u_t$ using its Euclidean divergence as (see Proposition 2.46 of \citealt{leeIntroductionRiemannianManifolds2018}):
\begin{equation*}
  \dive u_t(\Xm) = \dive_{\R} u_t(\Xm) - \frac{d+1}{2} \langle u_t(\Xm), \nabla_{\R} \log \det (\Xm)\rangle_{\R}%
\end{equation*}
where $\dive_{\R}$ is the Euclidean divergence. Finally, as done in \cite{lipmanFlowMatchingGenerative2023}, we use the unbiased Hutchinson trace estimator \citep{hutchinson1989stochastic} to estimate the Euclidean divergence $\dive_{\R} u_t$.

\section{More details on the synthetic experiments}
\label{appendix:details_synthetic_experiments}

In this section, we give more details on the synthetic experiments we led in \cref{sec:synthetic_exp}. We start by detailing the datasets and then give the setup and the architecture used.

\subsection{Details on the datasets}
\paragraph{The SPD Swiss-Roll}
We give here the equations that were used to generate the Swiss-roll dataset. We start by generating a 3D spiral using the following equations:
\begin{equation}
  \begin{aligned}
    u &= t\cos(t) \\
    v &= t\sin(t) \\
    w &= t + s + 1,
  \end{aligned}
\end{equation}
where $t \sim \mathcal{U}[0, 2\pi]$ is the angle of the spiral and $s \sim \mathcal{U}[0, 3]$ is the ribbon's width. Then, we rotate the 3D spiral to put it into the cone of equation $ac > b^2$ i.e. the cone of $2 \times 2$ SPD matrices of the form $
\begin{pmatrix}a&b\\b&c
\end{pmatrix}$:
\begin{equation}
  \begin{aligned}
    a &= \frac{1}{2}(w+u) \\
    c &= \frac{1}{2}(w-u) \\
    b &= \frac{1}{2}v.
  \end{aligned}
\end{equation}
Then, the matrix $$\Xm =
\begin{pmatrix}a&b\\b&c
\end{pmatrix}$$
will be SPD and will form a spiral in the cone of $2 \times 2$ SPD matrices.
\paragraph{The Figure-8}
Now, let us give the equations used to create the SPD Figure-8 dataset. We use a similar technic than for the above Swiss-Roll dataset. We start by creating a lemniscate-like shape using the following equations:
\begin{equation}
  \begin{aligned}
    u &= \cos(t) \\
    v &= \sin(2t) \\
    w &= t + s + \frac{1}{2},
  \end{aligned}
\end{equation}
where $t \sim \mathcal{U}[0, 2\pi]$ is the angle of the spiral and $s \sim \mathcal{U}[0, 1]$ is the ribbon's width. Then, we rotate the 3D spiral to put it into the cone of equation $ac > b^2$ i.e. the cone of $2 \times 2$ SPD matrices of the form $
\begin{pmatrix}a&b\\b&c
\end{pmatrix}$:
\begin{equation}
  \begin{aligned}
    a &= \frac{1}{2}(w+u) \\
    c &= \frac{1}{2}(w-u) \\
    b &= \frac{1}{2}v.
  \end{aligned}
\end{equation}
Then, the matrix $$\Xm =
\begin{pmatrix}a&b\\b&c
\end{pmatrix}$$
will be SPD and will form a 8 shape in the cone of $2 \times 2$ SPD matrices.

\paragraph{The Mixture of WG}
The last dataset is composed of $10 \times 10$ SPD matrices sampled from a mixture of three wrapped Gaussians. The mixture is:
$$q = \frac{1}{3}\WG(\Xm_1, \mu_1, \Sigma_1) + \frac{1}{3}\WG(\Xm_2, \mu_2, \Sigma_2) + \frac{1}{3}\WG(\Xm_3, \mu_3, \Sigma_3).$$
The center $\Xm_k$, for $k = 1,2,3$ are randomly generated as follows: we sample eigen-values $(\lambda_i)_{i = 1,...,d}$ uniformly in $[0.1, 10]$ and a random rotation matrix by sampling a random matrix and computing its QR decomposition. The $Q$ matrix we get is a orthogonal matrix. Finally, to create $\Xm_k$, we recompose $\Xm_k = Q \text{diag}(\lambda_1,...,\lambda_d)Q^\top$. The covariance matrices $\Sigma_k$ are sampled in the same way, with eigen-values sampled uniformly in $[0.05, 0.5]$. Finally,  $\mu_k$ is a random vector sampled uniformly in $[0,1]^d$. Finally, the whole dataset is reduced to have a variance of $3$, avoiding matrices with too big coefficients. We sample the same amount of points from each of the three wrapped Gaussian, and label them according to the component of the mixture we sampled from.

For all of those three datasets, the training dataset is composed of $8000$ matrices, the validation set is composed of $1000$ matrices and the testing set is composed of $1000$ matrices. We created 5 different versions of each datasets, each sampled with a seed ranging from $0$ et $4$.

\subsection{Details on the architecture used}
For the two fist $2 \times 2$ datasets (the \emph{SPD Swiss-roll} and the \emph{Figure-8}), the conditional vector field $u_t(\cdot \mid \Xm_1)$ was learned using a MLP composed of 6 layers, each composed of 128 neurons a Swish activation function \citep{ramachandran2017searchingactivationfunctions} with a learnable parameter. Timestep and class information are injected through a FiLM layer (\cite{perez2018film}). The input $\Xm$ of conditional vector field $u_t(\Xm \mid \Xm_1)$ is an SPD matrix and the output being a element of the tangent space $T_{\Xm} \P_d$, they are both symmetric matrices. Therefore, we vectorized the input matrix by taking only the elements on the superior diagonal and unvectorize the output matrix.

For the higher dimensional dataset, we used a more complex architecture composed of 12 layers each with 512 hidden layers a Swish activation function \citep{ramachandran2017searchingactivationfunctions} with a learnable parameter.

For all the synthetic experiments, we optimized the model using the Adam optimizer \citep{kingma2017adammethodstochasticoptimization} with a learning rate of $10^{-3}$ and and an exponential moving averaging on the weights \citep{PolyakAcceleration} with a decay of $10^{-5}$. We trained it for $50000$ epochs, with an early stopping of $100$ epochs patience triggered if the validation loss failed to improve. The synthetic experiments were conducted on NVIDIA A40 GPU with 40GB of memory. We train our Riemannian models with the AIRM loss used in \citet{chen2024flow} and the Euclidean CFM using the classical Frobenius loss on symmetric matrices.

\section{More results on the synthetics experiments}
\label{appendix:more_results_synthetics}
In this section, we give more quantitative metrics on the synthetic experiments we described in \cref{sec:synthetic_exp}. In addition to the metrics already considered in the core of the paper (the \% output matrices that are SPD, the log-likelihood and the Wasserstein-2 distance), we also measure the $\alpha$-Precision and $\beta$-Recall. We start by describing these metrics and then give the results on the synthetic datasets.

\paragraph{$\alpha$-Precision and $\beta$-Recall}
These $\alpha$-Precision and $\beta$-Recall are two metrics have been used for generative models used in image generation \citep{sajjadi2018assessing, kynkaanniemi2019improved} and have been adapted to LLMs \citep{bronnec2024exploring}. The $\alpha$-precision measures the quality of the generated samples (i.e. how realistic the generated samples are realistic or plausible compared to the true data distribution) while the $\beta$-recall quantifies the diversity of the generated samples (i.e. how well the true data distribution that is covered or represented by the generated samples). Both metrics range from 0 to 1, where higher values indicate better performance. We also report the harmonic mean of the two, denoted $\alpha, \beta$-F1 in order to balance fidelity and diversity into single metric. We used the algorithm provided in \citet{bronnec2024exploring} and adapted their code to the Riemannian case of SPD matrices.

\paragraph{The results}
We give the results of the $\alpha$-Precision and the $\beta$-Recall at \cref{table:res_synthetic_precision_recall}. We can see that the proposed WrappedCFM always has the best $\beta$-Recall, meaning that our method is the best at generating diverse samples. This is very important as we want to make sure that the generated samples cover a vast area of the distribution, and do not stay only in certain parts of the latter. However, our method does not achieve top performance in $\alpha$-Precision. The performance are still good, but are not enough to beat other methods.
\input{ICML/results_synthetic_exp_precision_recall}

\section{Details on the training and sampling times}
\label{appendix:training_times}

\paragraph{Training times}
\begin{table}[ht]
  \centering
  \small
  \setlength{\tabcolsep}{8pt}
  \begin{tabular}{lccc}
    \toprule
    \textbf{Model}
    & \textbf{Swiss Roll (s)}
    & \textbf{Figure 8 (s)}
    & \textbf{Mixture of WG}\\
    \midrule
    \textbf{WrappedCFM (proposed)}
    & $2795 \pm 760$
    & $3199 \pm 672$
    &$18652.32 \pm 2854.32$ \\

    Euclidean CFM
    & $2245 \pm 702$
    & $2471 \pm 361$
    & $2534.53 \pm 333.29$ \\

    DiffeoCFM
    & $\mathbf{1708} \pm 477$
    & $\mathbf{1813} \pm 502$
    & $7522.80 \pm 756.96$\\

    RiemannianCFM
    & $2129 \pm 494$
    & $2164 \pm 726$
    & $9919.13 \pm 1674.05$\\
    \bottomrule
  \end{tabular}
  \caption{Training time (mean $\pm$ standard deviation, in seconds) for different CFM variants on Swiss Roll and Figure~8 datasets.}
  \label{tab:training_time_synthetic_exp}
\end{table}

In this section, we report on the training time of the different methods on the different datasets. For the synthetic experiments, we give the results at \cref{tab:training_time_synthetic_exp}. We can see that our proposed method (WrappedCFM) is the slowest to train, which is coherent as it the conditional vector field $u_t(\cdot \mid \Xm_1)$ is the most complicated to compute (see \cref{eq:conditional_vector_field}) and it is computed at each iteration of the training.

\paragraph{Sampling times}

\begin{table}[ht]
  \centering
  \small
  \setlength{\tabcolsep}{8pt}
  \begin{tabular}{lccc}
    \toprule
    \textbf{Model}
    & \textbf{Swiss Roll (s)}
    & \textbf{Figure 8 (s)}
    & \textbf{Mixture of WG}\\
    \midrule
    WrappedCFM
    & $1.91 \pm 0.32$
    & $1.90 \pm 0.27$
    &$1.34 \pm 0.30$\\

    Euclidean CFM
    & $1.66 \pm 0.31$
    & $1.75 \pm 0.10$
    &$0.83     \pm  0.42$\\

    DiffeoCFM
    & $\mathbf{1.48} \pm 0.29$
    & $1.68 \pm 0.11$
    &$\mathbf{0.63}    \pm   0.06$\\

    RiemannianCFM
    & $1.69 \pm 0.32$
    & $\mathbf{1.66} \pm 0.42$
    &$1.96     \pm  1.00$\\
    \bottomrule
  \end{tabular}
  \caption{Sampling time (mean $\pm$ standard deviation, in seconds) for $1000$ samples using the different CFM variants on Swiss Roll and Figure~8 datasets.}
  \label{tab:sampling_time}
\end{table}
We also report the sampling times for $1000$ samples generating using the different methods at \cref{tab:sampling_time}. We can make the same comment as for the training time: the proposed WrappedCFM is the slowest to sample as it relies on an Riemannian Euler schema, as described at \cref{eq:euler_scheme}. Nevertheless, all methods have very reasonable sampling times, including WrappedCFM.
\section{More details on real-world experiments}

\subsection{Datasets preprocessing}
\label{appendix:datasets_preprocessing}

\begin{table}[ht]
  \centering
  \small
  \setlength{\tabcolsep}{6pt}
  \begin{tabular}{l l c c c l}
    \toprule
    \textbf{Dataset}
    & \textbf{Application}
    & \textbf{Dim.}
    & \textbf{\# Matrices (train/test)}
    & \textbf{\# Classes}
    & \textbf{Reference} \\
    \midrule
    BNCI2014004
    & BCI
    & $3 \times 3$
    & $5868 / 652$
    & $2$
    & \citep{Leebbcidataset} \\

    BNCI2014002
    & BCI
    & $15 \times 15$
    & $1{,}587 / 177$
    & $2$
    & \citep{SteyrlSchererFallerMüllerPutz+2016+77+86} \\

    BNCI2015001
    & BCI
    & $13 \times 13$
    & $3888 / 432$
    & $2$
    & \cite{faller} \\

    Textile
    & Image Analysis
    & $10 \times 10$
    & $14{,}400 / 1{,}600$
    & $2$
    & \cite{MVTec2021} \\
    \bottomrule
  \end{tabular}
  \caption{Summary of datasets used in the experiments.}
  \label{tab:datasets_summary}
\end{table}
% preprocessing of datasets
\subsection*{BCI}
The BCI datasets that we consider are \emph{BNCI2014004} \citep{Leebbcidataset}, \emph{BNCI2014002} \citep{SteyrlSchererFallerMüllerPutz+2016+77+86} and \emph{BNCI2015001} \citep{faller}. They consist of several subjects and several sessions per subjects performing a Motor Imagery task \cite{motor_imagery}. We used the library MOABB \cite{Aristimunha_Mother_of_all} to load and preprocess the data. The preprocess of BNCI2014004 consists of a band-pass filter in the range $[7;35]$ Hz. Then, we used the Ledoit-Wolf shrunk covariance matrix \citep{ledoitWellconditionedEstimatorLargedimensional2004} to in order to compute the covariance matrices and to avoid ill-conditioned matrices. For BNCI2014002 and BNCI2015001, we used the OAS estimator \citep{chen2010shrinkage} to compute the covariance matrices. We followed the same preprocessing as in \citet{collasRiemannianFlowMatching2025}.

\subsection*{Textile}
This dataset is made of a set of real images from textile manufacturing that contain non-defective and defective woven textiles. These images come from the public MVTec Anomaly Detection dataset \cite{MVTec2021}. The same preprocessing as in \citet{smithDataAnalysisUsing2022} was applied to get the covariance matrices. The two classes correspond to defective and non-defective textiles.

To ensure numerical stability, datasets are centered by the Riemannian mean and rescaled to unit Fréchet variance under the affine-invariant metric.

\subsection{Training details}
\label{appendix:real_training_details}

\paragraph{Flow matching models} The datasets are split into train, validation, test with 80/10/10 ratio. All the models are trained using the same MLP architecture as the synthetic experiments, with 12 layers and 512 hidden units. We trained until 100k iterations and used Adam optimizer with a learning rate of 2e-4 and exponential moving average on weights with a decay of 0.999 for all our experiments. All the experiments are repeated 5 times with different random seeds for data splits.

\paragraph{Classifier}

The classifier is the Riemannian Residual Network (RResNet) for SPD matrices using the log-euclidean metric \citet{katsman2023riemannianresnet} implemented by \citet{cov_logeucresnet_pelletier2025}, with 1 residual block, and 80\% of the original matrix size as the hyperparameter for the reduced dimensional space.

%All experiments are run on NVIDIA 2080 GPU with 11GB of memory. The network is a standard multi-layer perceptron taken from \cite{lipman2024flowmatchingguidecode} with Swish activation functions with a learnable parameter. For class conditional generation, we injected timestep and class information through a FiLM layer (\cite{perez2018film}) as done in U-nets used for image generation with diffusion models (\cite{diffusion_beats_gan}). We used Adam optimizer with a learning rate of 2e-4 and exponential moving average on weights with a decay of 0.999 for all our experiments.

\subsection{Additional metrics}
\input{ICML/results_datasets_app}
% \newpage
% % =====================================================================
% %  Geometric comparison of WrappedCFM vs DiffeoCFM vs RiemannianCFM
% %  Drop-in section/appendix. \input{wcfm_comparison} after your preamble.
% %
% %  Assumes your existing macros: \Xm \Ym \Im \Sm \Zm \Um \Vm \P \S \R \N
% %  \WG \Exp \Log \tr \dd  (from math_commands.tex / your preamble) and the
% %  amsthm environments theorem/proposition/lemma/corollary/remark/definition.
% %
% %  New helper commands are \providecommand-guarded to avoid clashes.
% % =====================================================================

% \providecommand{\sinhc}{\operatorname{sinhc}}
% \providecommand{\Mt}{\bm{M}_t}
% \providecommand{\cond}{\operatorname{cond}}
% \providecommand{\dexp}{\mathrm{d}\!\exp}
% \providecommand{\iotaX}[1]{\iota_{#1}}
% \providecommand{\siginit}{\sigma_{\mathrm{init}}}
% \providecommand{\sigmin}{\sigma_{\min}}

% \section{Why WrappedCFM? A geometric comparison}
% \label{sec:comparison}

% In this section we compare, at the level of the conditional probability
% paths and their flows, the three members of our framework: RiemannianCFM,
% DiffeoCFM and the proposed WrappedCFM. The guiding observation is the
% following. WrappedCFM encodes the \emph{location} of the conditional path
% \emph{intrinsically}, through a base point $\bar{\Xm}_t=\Xm_1^{t}$ that rides
% the geodesic, while DiffeoCFM encodes it \emph{extrinsically}, through a
% fixed chart anchored at $\Im_d$ together with a tangent--space mean
% $\mu_t=t\log\Xm_1$. A single chart cannot keep the affine--invariant metric
% isotropic away from its base point, and we show that the resulting
% distortion is \emph{exactly} the $\sinh$--Jacobian of the exponential map --
% the same object that appears in \cref{prop:jac_exp}. RiemannianCFM is the
% degenerate, zero--noise corner of the same family.

% \paragraph{Conventions.}
% We use the affine--invariant (AIRM--orthonormal) identification of each
% tangent space with $\R^D$: a vector $\Um\in T_{\Xm}\P_d$ is represented by
% $\iotaX{\Xm}(\Um):=\Xm^{-1/2}\Um\Xm^{-1/2}$, vectorized in a fixed
% Frobenius--orthonormal basis of $\S_d$, so that
% $\langle\Um,\Vm\rangle_{\Xm}=\langle\iotaX{\Xm}\Um,\iotaX{\Xm}\Vm\rangle_F$.
% Throughout, $c_t:=\gamma_{\Im_d,\Xm_1}(t)=\Xm_1^{t}$ denotes the
% constant--speed AIRM geodesic from $\Im_d$ to $\Xm_1$, and
% $\sinhc(x):=\sinh(x)/x$ (with $\sinhc(0)=1$). We recall the linear scheduler
% $\sigma_t=\siginit-(\siginit-\sigmin)t$, so that $\sigma_t>0$ and
% $\sigma_t'=-(\siginit-\sigmin)<0$.

% Both WrappedCFM and DiffeoCFM center their conditional law on the
% \emph{same} geodesic point $c_t$: for WrappedCFM the base point is $\Xm_1^t$
% with zero mean, and for DiffeoCFM the base point is $\Im_d$ with mean
% $t\log\Xm_1$, which yields $\Exp_{\Im_d}(t\log\Xm_1)=\Xm_1^{t}=c_t$. Hence
% the mean trajectories coincide, and the entire difference between the two
% methods lies in the \emph{covariance} of the conditional law. This is what we
% now quantify.

% % ---------------------------------------------------------------------
% \subsection{WrappedCFM is exactly isotropic in the moving normal frame}
% % ---------------------------------------------------------------------

% \begin{proposition}[Exact isotropy]
% \label{prop:wcfm_isotropy}
% Let $\Xm\sim p_0=\WG(\Im_d,0_D,\siginit^2\Im_D)$ and
% $\Ym_t=\Phi^{\mathrm{WCFM}}_t(\Xm)=\Xm_1^{t/2}\Xm^{\sigma_t/\siginit}\Xm_1^{t/2}$.
% Then, in AIRM--orthonormal coordinates at the moving center $c_t=\Xm_1^{t}$,
% \begin{equation}
% \label{eq:wcfm_normal_coord}
% \iotaX{c_t}\!\big(\Log_{c_t}\Ym_t\big)
% =\frac{\sigma_t}{\siginit}\,\log\Xm
% \;\sim\;\mathcal N\!\big(0_D,\sigma_t^2\Im_D\big).
% \end{equation}
% Consequently $p^{\mathrm{WCFM}}_t(\cdot\mid\Xm_1)$ is, in normal coordinates
% at $c_t$, \emph{exactly} the isotropic Gaussian $\mathcal N(0_D,\sigma_t^2\Im_D)$
% -- condition number $1$ for every $\Xm_1$ and every $t$ -- and the flow acts
% in that frame as the pure isotropic scaling
% $\mathbf w\mapsto(\sigma_t/\siginit)\,\mathbf w$.
% \end{proposition}

% \begin{proof}
% A direct computation gives
% $c_t^{-1/2}\Ym_t c_t^{-1/2}
% =\Xm_1^{-t/2}\big(\Xm_1^{t/2}\Xm^{\sigma_t/\siginit}\Xm_1^{t/2}\big)\Xm_1^{-t/2}
% =\Xm^{\sigma_t/\siginit}$,
% so that
% $\log\!\big(c_t^{-1/2}\Ym_t c_t^{-1/2}\big)=\tfrac{\sigma_t}{\siginit}\log\Xm$
% and therefore
% $\iotaX{c_t}(\Log_{c_t}\Ym_t)
% =c_t^{-1/2}\Log_{c_t}(\Ym_t)c_t^{-1/2}
% =\tfrac{\sigma_t}{\siginit}\log\Xm$.
% Since $\Xm\sim\WG(\Im_d,0_D,\siginit^2\Im_D)$ means
% $\log\Xm\sim\mathcal N(0_D,\siginit^2\Im_D)$, the right-hand side is
% $\mathcal N(0_D,\sigma_t^2\Im_D)$. The statement is exact because $\P_d$ is a
% Hadamard manifold, so $\Log_{c_t}\circ\Exp_{c_t}=\mathrm{id}$ on all of
% $\S_d$.
% \end{proof}

% Equation~\eqref{eq:wcfm_normal_coord} shows that the conditional cloud never
% rotates or shears: a single realization $\log\Xm$ is merely rescaled across
% time. This is the precise sense in which WrappedCFM is geometrically
% faithful.

% % ---------------------------------------------------------------------
% \subsection{DiffeoCFM's anisotropy is the exponential-map Jacobian}
% % ---------------------------------------------------------------------

% Let $\bm L:=t\log\Xm_1=\sum_i\ell_i\,\bm u_i\bm u_i^\top$ with
% $\ell_i=t\lambda_i$, where $\lambda_1,\dots,\lambda_d$ are the eigenvalues of
% $\log\Xm_1$. Let $\{\bm E_{ij}\}_{i\le j}$ be the induced
% Frobenius--orthonormal eigenbasis of $\S_d$, i.e.\ $\bm E_{ii}=\bm u_i\bm u_i^\top$
% and $\bm E_{ij}=\tfrac1{\sqrt2}(\bm u_i\bm u_j^\top+\bm u_j\bm u_i^\top)$ for
% $i<j$.

% \begin{proposition}[Anisotropy spectrum of DiffeoCFM]
% \label{prop:diffeo_anisotropy}
% Consider the DiffeoCFM noise--injection map
% $\bm N^{\mathrm D}_t:\xi\mapsto\Exp_{\Im_d}(\bm L+\xi)$, with
% $\bm N^{\mathrm D}_t(0)=c_t$. Its differential at $0$, viewed as a linear map
% $(T_{\Im_d}\P_d,\langle\cdot\rangle_{\Im_d})\to(T_{c_t}\P_d,\langle\cdot\rangle_{c_t})$,
% is diagonal in $\{\bm E_{ij}\}$ with singular values
% \begin{equation}
% \label{eq:tau_ij}
% \tau_{ij}=\sinhc\!\Big(\tfrac{\ell_i-\ell_j}{2}\Big)
% =\sinhc\!\Big(\tfrac{t(\lambda_i-\lambda_j)}{2}\Big)\;\ge\;1,
% \qquad \tau_{ii}=1 .
% \end{equation}
% Hence the DiffeoCFM conditional covariance in normal coordinates at $c_t$ is,
% to leading order in $\sigma_t$, equal to $\sigma_t^2\operatorname{diag}(\tau_{ij}^2)$,
% whereas WrappedCFM's is $\sigma_t^2\Im_D$ exactly. Writing
% $\Delta=\lambda_{\max}-\lambda_{\min}$ for the spectral spread of $\log\Xm_1$,
% \begin{equation}
% \label{eq:cond_and_det}
% \kappa\big(\Sigma^{\mathrm D}_t\big)=\sinhc\!\Big(\tfrac{t\Delta}{2}\Big)^{2},
% \qquad
% \frac{\det\Sigma^{\mathrm D}_t}{\det\Sigma^{\mathrm W}_t}
% =\prod_{i<j}\tau_{ij}^2
% =J_{\Im_d}\!\big(t\log\Xm_1\big)^2 .
% \end{equation}
% \end{proposition}

% \begin{proof}
% Using $\dexp_{\bm L}(\eta)=\int_0^1 e^{s\bm L}\eta\,e^{(1-s)\bm L}\,ds$,
% the matrix $\dexp_{\bm L}$ is diagonal on $\{\bm E_{ij}\}$ with eigenvalue
% $\int_0^1 e^{s\ell_i}e^{(1-s)\ell_j}ds=\tfrac{e^{\ell_i}-e^{\ell_j}}{\ell_i-\ell_j}$.
% The differential of $\bm N^{\mathrm D}_t$ at $0$ is $\dexp_{\bm L}$, and its
% AIRM--orthonormal coordinate at $c_t=e^{\bm L}$ is
% \[
% \iotaX{c_t}\big(\dexp_{\bm L}(\bm E_{ij})\big)
% =e^{-\bm L/2}\,\frac{e^{\ell_i}-e^{\ell_j}}{\ell_i-\ell_j}\,\bm E_{ij}\,e^{-\bm L/2}
% =e^{-(\ell_i+\ell_j)/2}\,\frac{e^{\ell_i}-e^{\ell_j}}{\ell_i-\ell_j}\,\bm E_{ij}
% =\tau_{ij}\,\bm E_{ij},
% \]
% because $\frac{e^{a}-e^{b}}{a-b}e^{-(a+b)/2}
% =\frac{2\sinh((a-b)/2)}{a-b}=\sinhc(\tfrac{a-b}{2})$. As $\bm E_{ij}$ has unit
% AIRM norm at $\Im_d$, the $\tau_{ij}$ are precisely the singular values.
% Since $\sinhc\ge1$ and is even and increasing in $|x|$, the smallest singular
% value is $1$ (attained on the diagonal modes) and the largest is
% $\sinhc(t\Delta/2)$, which gives the condition number in
% \eqref{eq:cond_and_det}. The leading--order covariance follows from the
% normal--coordinate (Laplace) expansion
% $\iotaX{c_t}(\Log_{c_t}\Ym)=\sigma_t\,\dexp_{\bm L}(\xi)\,/$ frame $+O(\sigma_t^2)$
% with $\Ym=\Exp_{\Im_d}(\bm L+\sigma_t\xi)$, $\xi\sim\mathcal N(0,\Im_D)$.
% Finally
% $\prod_{i<j}\tau_{ij}
% =\prod_{i<j}2\,\tfrac{\sinh((\ell_i-\ell_j)/2)}{\ell_i-\ell_j}
% =2^{d(d-1)/2}\prod_{i<j}\tfrac{\sinh((\ell_i-\ell_j)/2)}{\ell_i-\ell_j}
% =J_{\Im_d}(\bm L)$ by \cref{prop:jac_exp}, and the diagonal modes contribute
% $\tau_{ii}=1$, so the determinant ratio is $J_{\Im_d}(t\log\Xm_1)^2$.
% \end{proof}

% \begin{corollary}[The condition number tracks $\cond(\Xm_1)$]
% \label{cor:cond}
% As $\Delta\to\infty$ (with $t>0$ fixed),
% $\kappa(\Sigma^{\mathrm D}_t)\sim e^{t\Delta}/(t\Delta)^2$. In particular at
% $t=1$, since $\Delta=\log\cond(\Xm_1)$,
% \[
% \kappa\big(\Sigma^{\mathrm D}_1\big)\;\sim\;
% \frac{\cond(\Xm_1)}{\big(\log\cond(\Xm_1)\big)^2}.
% \]
% The DiffeoCFM conditional anisotropy therefore grows essentially linearly
% with the condition number of the target matrix $\Xm_1$ -- exactly the
% ill--conditioned regime typical of EEG and DTI covariance matrices -- while
% WrappedCFM remains isotropic at $\kappa\equiv1$.
% \end{corollary}

% \begin{remark}[Mis-shapen, not less diverse]
% Because $\tau_{ij}\ge1$, DiffeoCFM injects \emph{more} total noise volume
% than WrappedCFM ($\det\Sigma^{\mathrm D}_t\ge\det\Sigma^{\mathrm W}_t$): it
% over--disperses the ``mixing'' off--diagonal modes while leaving the
% diagonal modes at scale $\sigma_t$. The relevant defect is thus the
% \emph{shape} of the conditional noise, measured by the scale--free condition
% number, rather than its overall magnitude.
% \end{remark}

% % ---------------------------------------------------------------------
% \subsection{The WrappedCFM field: geodesic translation plus isotropic contraction}
% % ---------------------------------------------------------------------

% \begin{proposition}[Structure of the conditional field]
% \label{prop:field_split}
% The WrappedCFM conditional vector field decomposes as
% \[
% u^{\mathrm{WCFM}}_t(\Ym\mid\Xm_1)
% =\underbrace{\tfrac12\big(\log(\Xm_1)\,\Ym+\Ym\,\log(\Xm_1)\big)}_{=:~\mathcal T(\Ym)}
% \;+\;
% \underbrace{\frac{\sigma_t'}{\sigma_t}\,\Xm_1^{t/2}\Mt^{1/2}\log(\Mt)\Mt^{1/2}\Xm_1^{t/2}}_{=:~(\sigma_t'/\sigma_t)\,\mathcal C_t(\Ym)},
% \quad \Mt=\Xm_1^{-t/2}\Ym\,\Xm_1^{-t/2},
% \]
% with the following properties:
% \begin{enumerate}
% \item[(i)] $\mathcal T$ is the infinitesimal generator of the
% geodesic--translation isometry $\Ym\mapsto\Xm_1^{s/2}\Ym\,\Xm_1^{s/2}$;
% \item[(ii)] at the center $\Ym=c_t$ the contraction vanishes and
% $u^{\mathrm{WCFM}}_t(c_t\mid\Xm_1)=\dot\gamma_{\Im_d,\Xm_1}(t)$, the exact
% geodesic velocity;
% \item[(iii)] in the moving normal frame the contraction is exactly isotropic:
% the coordinate $\mathbf w_t=\iotaX{c_t}(\Log_{c_t}\Ym_t)$ obeys
% $\dot{\mathbf w}_t=\tfrac{\sigma_t'}{\sigma_t}\,\mathbf w_t$ with
% $\sigma_t'/\sigma_t<0$ (a contraction toward the geodesic).
% \end{enumerate}
% \end{proposition}

% \begin{proof}
% (i) $\tfrac{\dd}{\dd s}\big|_{s=0}\Xm_1^{s/2}\Ym\,\Xm_1^{s/2}
% =\tfrac12\log(\Xm_1)\Ym+\tfrac12\Ym\log(\Xm_1)=\mathcal T(\Ym)$, and
% $\Ym\mapsto\Xm_1^{s/2}\Ym\,\Xm_1^{s/2}$ is a congruence by an SPD matrix,
% hence an AIRM isometry. (ii) At $\Ym=c_t=\Xm_1^t$ we have
% $\Mt=\Xm_1^{-t/2}\Xm_1^t\Xm_1^{-t/2}=\Im_d$, so $\log\Mt=0$ and the
% contraction vanishes; moreover $\mathcal T(c_t)=\log(\Xm_1)\Xm_1^{t}
% =\tfrac{\dd}{\dd t}\Xm_1^{t}=\dot\gamma_{\Im_d,\Xm_1}(t)$, using that
% $\log\Xm_1$ and $\Xm_1^{t}$ commute. (iii) By \cref{prop:wcfm_isotropy},
% $\mathbf w_t=\tfrac{\sigma_t}{\siginit}\log\Xm$, hence
% $\dot{\mathbf w}_t=\tfrac{\sigma_t'}{\siginit}\log\Xm
% =\tfrac{\sigma_t'}{\sigma_t}\mathbf w_t$.
% \end{proof}

% Thus WrappedCFM is, in one phrase, ``the RiemannianCFM geodesic field plus a
% manifold--aware restoring drift that anneals the conditional cloud
% isotropically from $\siginit$ to $\sigmin$.'' DiffeoCFM performs the same
% annealing but inside a single tangent space, which is exactly why its cloud
% shears according to \cref{prop:diffeo_anisotropy}.

% % ---------------------------------------------------------------------
% \subsection{WrappedCFM strictly generalizes RiemannianCFM}
% % ---------------------------------------------------------------------

% \begin{proposition}[Zero--noise limit]
% \label{prop:rcfm_limit}
% Scale the schedule as $\sigma_t^{(\varepsilon)}=\varepsilon\,\bar\sigma_t$ with
% $\bar\sigma_t=a-(a-b)t$, $a,b>0$. As $\varepsilon\to0$:
% \begin{enumerate}
% \item[(i)] $p^{(\varepsilon)}_t(\cdot\mid\Xm_1)\to\delta_{\gamma_{\Im_d,\Xm_1}(t)}$
% weakly, for each $t\in[0,1]$;
% \item[(ii)] on the (shrinking) support of the conditional law,
% $u^{(\varepsilon)}_t(\cdot\mid\Xm_1)\to\dot\gamma_{\Im_d,\Xm_1}(t)$.
% \end{enumerate}
% Hence RiemannianCFM is exactly the degenerate $\sigma_t\equiv0$ corner of the
% framework, and WrappedCFM is its full--support deformation.
% \end{proposition}

% \begin{proof}
% (i) The conditional law is a wrapped Gaussian with covariance
% $\varepsilon^2\bar\sigma_t^2\Im_D\to0$, hence concentrates at its center
% $c_t=\gamma_{\Im_d,\Xm_1}(t)$. (ii) The log--derivative
% $\sigma_t'/\sigma_t=\bar\sigma_t'/\bar\sigma_t$ is scale invariant and $O(1)$.
% On the support, $\Ym=\Exp_{c_t}(O(\varepsilon))$ gives
% $\Mt=\Im_d+O(\varepsilon)$, so $\log\Mt=O(\varepsilon)$ and the contraction
% term is $O(\varepsilon)\to0$, leaving $\mathcal T(c_t)=\dot\gamma_{\Im_d,\Xm_1}(t)$
% by \cref{prop:field_split}(ii).
% \end{proof}

% \begin{remark}[Denser supervision]
% With $\sigma_t>0$ the regression target $u_t(\cdot\mid\Xm_1)$ is defined on
% all of $\P_d$ at each $t$, so the marginal field is supervised on a
% full--dimensional support; RiemannianCFM ($\sigma_t\equiv0$) only supervises
% on the measure--zero union of geodesics. The trade--off is that RiemannianCFM
% reaches $q_1$ exactly at $t=1$ (its terminal conditional is $\delta_{\Xm_1}$),
% whereas WrappedCFM and DiffeoCFM recover a $\sigmin$--smoothed $q_1$.
% \end{remark}

% % ---------------------------------------------------------------------
% \subsection{Equivariance of the framework}
% % ---------------------------------------------------------------------

% \begin{proposition}[$O(d)$-- and inversion--equivariance, and maximality]
% \label{prop:equivariance}
% For every $\bm O\in O(d)$,
% $\Phi^{\bm O\Xm_1\bm O^\top}_t(\bm O\Xm\bm O^\top)=\bm O\,\Phi^{\Xm_1}_t(\Xm)\,\bm O^\top$,
% and the source $\WG(\Im_d,0_D,\siginit^2\Im_D)$ is $O(d)$--invariant; hence the
% WrappedCFM generative model is $O(d)$--equivariant. It is moreover equivariant
% under inversion $\Xm\mapsto\Xm^{-1}$. No method whose source is anchored at
% $\Im_d$ can be equivariant under the full congruence group $GL_d$, so $O(d)$
% is the maximal congruence symmetry attainable.
% \end{proposition}

% \begin{proof}
% Orthogonal congruence commutes with analytic matrix functions:
% $(\bm O\Xm_1\bm O^\top)^{t/2}=\bm O\Xm_1^{t/2}\bm O^\top$ and
% $(\bm O\Xm\bm O^\top)^{\sigma_t/\siginit}=\bm O\Xm^{\sigma_t/\siginit}\bm O^\top$,
% so the three factors of $\Phi_t$ transform by conjugation and recombine.
% The source is $O(d)$--invariant because $\Im_d$ is fixed by $O(d)$ and the
% AIRM--orthonormal noise is rotation invariant. For inversion,
% $\Xm_1^{-t/2}\Xm^{-\sigma_t/\siginit}\Xm_1^{-t/2}
% =\big(\Xm_1^{t/2}\Xm^{\sigma_t/\siginit}\Xm_1^{t/2}\big)^{-1}$.
% Maximality: under congruence by $g\in GL_d$, $\Im_d\mapsto gg^\top\neq\Im_d$
% unless $g\in O(d)$, so the identity--anchored source cannot be
% $GL_d$--invariant.
% \end{proof}

% DiffeoCFM enjoys the same symmetry; Euclidean CFM does not. We record it as a
% structural correctness property of the framework.

% % ---------------------------------------------------------------------
% \subsection{A canonical characterization}
% % ---------------------------------------------------------------------

% \begin{proposition}[WrappedCFM is the canonical AIRM--natural lift]
% \label{prop:characterization}
% Among conditional paths
% $p_t(\cdot\mid\Xm_1)=\WG(\bar{\Xm}_t,\mu_t,\Sigma_t)$ conditioned on a single
% $\Xm_1\sim q_1$, with boundary conditions
% $\bar{\Xm}_0=\Im_d,\ \mu_0=0_D,\ \Sigma_0=\siginit^2\Im_D$ and terminal
% concentration at $\Xm_1$, WrappedCFM is the unique non--degenerate member
% that simultaneously satisfies: (a) the base point is the constant--speed
% geodesic, $\bar{\Xm}_t=\gamma_{\Im_d,\Xm_1}(t)$; (b) the tangent mean is
% zero, $\mu_t\equiv0_D$; and (c) the covariance is isotropic,
% $\Sigma_t\propto\Im_D$.
% \end{proposition}

% \begin{proof}[Sketch]
% DiffeoCFM keeps $\bar{\Xm}_t=\Im_d$ and carries all location in
% $\mu_t=t\log\Xm_1$, violating (a) and (b); RiemannianCFM sets $\Sigma_t=0$,
% violating (c) and lying on the degenerate boundary. WrappedCFM is the
% parameterization $\bar{\Xm}_t=\Xm_1^{t}$, $\mu_t=0_D$,
% $\Sigma_t=\sigma_t^2\Im_D$, which satisfies (a)--(c); the boundary conditions
% and isotropy with geodesic base point fix it uniquely.
% \end{proof}

% Condition (b) is the crux of the intrinsic--versus--extrinsic distinction:
% $\mu_t\equiv0$ means the conditional law is centered at its base point, so
% the path's location is carried by a \emph{manifold point} rather than by a
% tangent vector in a fixed chart. \Cref{prop:diffeo_anisotropy} shows that it
% is precisely the use of a single chart -- not the geometry of $\P_d$ -- that
% forces DiffeoCFM's anisotropy.

% % ---------------------------------------------------------------------
% \subsection{Figures}
% % ---------------------------------------------------------------------

% \begin{figure}[t]
%   \centering
%   \includegraphics[width=0.62\linewidth]{figures/fig_condition_number.pdf}
%   \caption{Condition number $\kappa_t=\sinhc(t\Delta/2)^2$ of the DiffeoCFM
%   conditional covariance at the geodesic point, as a function of the spectral
%   spread $\Delta=\log\cond(\Xm_1)$ of $\log\Xm_1$, for several $t$ (\cref{prop:diffeo_anisotropy,cor:cond}).
%   WrappedCFM is isotropic ($\kappa\equiv1$) for all $t$ and all targets.}
%   \label{fig:cond_number}
% \end{figure}

% \begin{figure}[t]
%   \centering
%   \includegraphics[width=0.85\linewidth]{figures/fig_ellipses.pdf}
%   \caption{Conditional noise in Riemannian normal coordinates along the
%   geodesic $\gamma_{\Im_d,\Xm_1}$ ($d{=}2$, $\Delta{=}8$). WrappedCFM (top)
%   stays an isotropic circle of radius $\sigma_t$ that anneals from $\siginit$
%   to $\sigmin$; DiffeoCFM (bottom) shears into an ellipse whose aspect ratio
%   $\tau_t=\sinhc(t\Delta/2)$ grows with $t$.}
%   \label{fig:ellipses}
% \end{figure}

% \begin{figure}[t]
%   \centering
%   \includegraphics[width=0.62\linewidth]{figures/fig_heatmap.pdf}
%   \caption{$\log_{10}\kappa_t$ for DiffeoCFM over time $t$ and target spread
%   $\Delta$. WrappedCFM corresponds to the identically--zero level.}
%   \label{fig:heatmap}
% \end{figure}

% \newpage

% \begin{theorem}
%     Let $p^{\text{WCFM}}_t(\cdot \mid \Xm_1)$ be the conditional probability path given of the WrappedCFM and $\Phi^{\text{WCFM}}_t$ its corresponding
% flow map as in \cref{eq:flow}. Then, the unique vector field that defines $\Phi^{\text{WCFM}}_t$ has the form:
% \begin{equation*}
% \begin{split}
%     u^{\text{WCFM}}_t(&\Ym \mid \Xm_1)
%   = \Sm(t)\bar{\Xm}_t^{-1/2}\,\Ym + \Ym \bar{\Xm}_t^{-1/2}\Sm(t)\\
%   &+ \tfrac{\sigma'_t(\Xm_1)}{\sigma_t(\Xm_1)} \,\bar{\Xm}_t^{1/2} \bm{M}_t^{1/2}\log(\bm{M}_t)\bm{M}_t^{1/2}\bar{\Xm}_t^{1/2},
% \end{split}
% \end{equation*}
% where $\Sm$ is the derivative of $t \mapsto \bar{\Xm}_t^{1/2}$ and $\bm{M}_t = \bar{\Xm}_t^{-1/2}\Ym \bar{\Xm}_t^{-1/2}$. Consequently, $u_t(\cdot \mid \Xm_1)$ generates the wrapped Gaussian path $p_t(\cdot \mid \Xm_1)$.
% \end{theorem}

\end{document}
